\chapter{Sums, Binomials and Induction}\label{ch:sums}

\begin{center}\itshape
Is it a coincidence that $1^3+2^3+\dots+n^3=(1+2+\dots+n)^2$?\\
By the end of this chapter you will have four reasons why it is not.
\end{center}

Chapter~\ref{ch:language} fixed the meaning of one kind of sentence: ``for every positive integer $n$''. A check of the first few values does not finish that sentence. It only checks those values. This chapter takes one such sentence and builds the tools that turn it into a proof.

Here is the sentence. For every integer $n\ge1$,
\[
1^3+2^3+\dots+n^3=(1+2+\dots+n)^2.
\]
The first four cases are easy to compute by hand. For $n=1$ both sides are $1$. For $n=2$, the left side is $1+8=9$ and the right side is $3^2$. For $n=3$, one more cube gives $9+27=36=6^2$. For $n=4$, one more gives $36+64=100=10^2$. The pattern is real for these four numbers. The sentence claims every larger $n$ as well.

Around the year $100$, Nicomachus of Gerasa recorded a grouping that produces the same numbers. He did not prove the identity. We prove it four times, after the tools are ready: by induction, by a change of basis, by his grouping of the odd numbers, and by cutting a multiplication table into L-shapes. All four are in Section~\ref{sec:fourproofs}.

A list of checks is not yet the sentence. The workshop shows a rule that matches the powers of two at five integers and fails at the sixth. When a pattern has worked for every integer you have tried, ask it about the next one.

The four proofs need a way to write a sum, a difference that makes a sum collapse, induction, and the numbers that count subsets. Those numbers, read as polynomials, sum every polynomial. Divisibility comes last, because the exercises on primes and Fibonacci numbers need it.

\section{Sums, and how to write them}\label{sec:sums}

A sum needs a notation, and three rules that cover almost every later move. The letter under the sigma sign is not a quantity you solve for. It is a place-holder for the terms.

\begin{reminder}[sigma notation]
For integers $m\le n$ and numbers $a_m,\dots,a_n$, write $\sum_{k=m}^na_k=a_m+a_{m+1}+\dots+a_n$. The letter $k$ is a \emph{dummy} index. It has no meaning outside the sum, and $\sum_{k=m}^na_k=\sum_{j=m}^na_j$. An empty sum, the case $m>n$, equals $0$. Three rules cover the usual moves.
\begin{enumerate}[label=(\roman*)]
\item \emph{Linearity}: $\sum(\lambda a_k+\mu b_k)=\lambda\sum a_k+\mu\sum b_k$.
\item \emph{Index shift}: $\sum_{k=m}^na_k=\sum_{j=m+1}^{n+1}a_{j-1}$, by putting $j=k+1$. More generally, the index may be renamed by any bijection of the range.
\item \emph{Splitting}: $\sum_{k=m}^na_k=\sum_{k=m}^pa_k+\sum_{k=p+1}^na_k$ whenever $m\le p<n$.
\end{enumerate}
A double sum $\sum_{i=1}^m\sum_{j=1}^na_{ij}$ adds the $mn$ entries of a rectangle. Rows first or columns first, the total is the same, so the two sigma signs may be swapped. When the inner range depends on the outer index, as in $\sum_{i=1}^n\sum_{j=1}^ia_{ij}$, swapping means describing the same set of pairs in the other order: $\sum_{j=1}^n\sum_{i=j}^na_{ij}$. Draw the pairs before you swap. Figure~\ref{fig:region} is the picture. Most wrong limits in this chapter come from skipping it.
\end{reminder}

The empty sum equals $0$ because $0$ is the number that leaves a total unchanged. Adding no terms should do nothing. An empty pocket adds nothing to your change.

\begin{figure}[ht]
\centering
\fig{ch01-region}
\caption{The same ten pairs $(i,j)$ with $1\le j\le i\le 4$. On the left, the outer sum fixes the column $i$ and the inner sum runs down that column. On the right, the outer sum fixes the row $j$ and the inner sum runs along that row. The shaded cells do not change. Only the order of addition changes. The pale cells are outside the sum.}
\label{fig:region}
\end{figure}

Check the empty sum and a shift on a sum you can do in your head. The sum $\sum_{k=5}^{4}k$ has no terms, so it is $0$. The sum $\sum_{k=1}^{3}k=6$ is the same list as $\sum_{j=2}^{4}(j-1)$, because $j=k+1$ runs from $2$ to $4$ while $k$ runs from $1$ to $3$.

Two formulas come up again and again. Pairing gives the first. A shift gives the second.

\begin{propositionx}[finite progressions]
Let $n\ge0$.
\begin{enumerate}[label=(\alph*)]
\item $\displaystyle\sum_{k=0}^nk=\frac{n(n+1)}{2}$. More generally, for any numbers $a$ and $d$,
\[
\sum_{k=0}^n(a+kd)=(n+1)a+d\,\frac{n(n+1)}{2}.
\]
\item If $q\ne1$, then $\displaystyle\sum_{k=0}^nq^k=\frac{q^{n+1}-1}{q-1}$. If $q=1$, the sum is $n+1$.
\end{enumerate}
\end{propositionx}
\begin{proof}
We prove (a) by pairing, then (b) by an index shift.

\emph{The sum of the first integers.} Let $S=\sum_{k=0}^nk$. The substitution $j=n-k$ reverses the list, and $j$ runs through the same range $0,1,\dots,n$. So $S=\sum_{j=0}^n(n-j)$. Rename the dummy index $j$ back to $k$, and add the two expressions for $S$ term by term:
\[
2S=\sum_{k=0}^n\bigl(k+(n-k)\bigr)=\sum_{k=0}^nn=(n+1)n.
\]
Divide by $2$. For $n=0$ both sides are $0$. For $n=6$, Figure~\ref{fig:gauss} is this identity: two copies of $1+\dots+6$ fill a $6\times7$ rectangle, so $2(1+\dots+6)=6\cdot7$.

\emph{The general arithmetic progression.} Split the sum by linearity:
\[
\sum_{k=0}^n(a+kd)=a\sum_{k=0}^n1+d\sum_{k=0}^nk=(n+1)a+d\,\frac{n(n+1)}{2}.
\]
There are $n+1$ terms in the first sum, because the index starts at $0$.

\emph{The geometric sum.} If $q=1$, then every term is $1$, and there are $n+1$ of them. Now suppose $q\ne1$. Let $G=\sum_{k=0}^nq^k$. Multiply by $q$ and shift the index by $j=k+1$:
\[
qG=\sum_{k=0}^nq^{k+1}=\sum_{j=1}^{n+1}q^j.
\]
Subtract. Every power $q^1,\dots,q^n$ appears in both $qG$ and $G$, so those terms cancel:
\[
qG-G=q^{n+1}-q^0=q^{n+1}-1.
\]
The left side is $(q-1)G$. Since $q\ne1$, we may divide:
\[
G=\frac{q^{n+1}-1}{q-1}.
\]
For a check, take $q=2$ and $n=3$. The sum is $1+2+4+8=15$, and $(2^4-1)/(2-1)=15$. Figure~\ref{fig:geometric} is that check drawn as a bar.
\end{proof}

\begin{figure}[ht]
\centering
\fig{ch01-geometric}
\caption{The blocks of length $1$, $2$, $4$ and $8$, in charcoal, coral, ochre and rust, fill fifteen units. The outlined unit at the end is the one not included. Together with it, the bar would have length $16=2^4$. So $1+2+4+8=2^4-1$, which is the geometric sum at $q=2$ and $n=3$.}
\label{fig:geometric}
\end{figure}

The arithmetic progression gives $\sum_{k=1}^nk=n(n+1)/2$, because the $k=0$ term is $0$. The geometric sum explains a one-line count you may already know: the binary number with eleven ones is one less than $2^{11}$, so $1+2+\dots+2^{10}=2^{11}-1$. That is the formula at $q=2$. The formula also covers every other ratio. That closed form is the one later geometric pieces use, from this binary count to the geometric series in Chapter~\ref{ch:sequences}.

There is a famous classroom story about the pairing. The earliest printed version is in the memorial that Wolfgang Sartorius von Waltershausen, a friend of the old Gauss, published in $1856$, the year after Gauss's death. It says that the task was an arithmetic progression, and that the slate landed on the table with a single number on it. The particular sum from $1$ to $100$, and the motive of keeping a class busy, belong to later retellings. Brian Hayes, in $2006$, collected more than a hundred versions. The pairing stands on its own. For the sum from $1$ to $100$, each of the hundred columns adds to $101$, so twice the sum is $100\cdot101$ and the sum is $5050$. Figure~\ref{fig:gauss} draws the same pairing for $n=6$. The argument works for every arithmetic progression.

\begin{figure}[ht]
\centering
\fig{ch01-gauss}
\caption{Left: two copies of $1+2+\dots+6$ fit as a $6\times7$ rectangle, so the sum is $6\cdot7/2$. Right: the odd numbers are L-shaped borders of a square of side $1+2+3+4=10$. Grouped as $1$, then $3+5$, then $7+9+11$, then $13+15+17+19$, those borders have areas $1^3,2^3,3^3,4^3$. Section~\ref{sec:fourproofs} turns the picture into a proof.}
\label{fig:gauss}
\end{figure}

\portrait{gauss}{Carl Friedrich Gauss (1777--1855). Portrait by Christian Albrecht Jensen, 1840. Public domain, via Wikimedia Commons.}

\begin{examplex}
What is $\sum_{k=1}^n(2k-1)$, the sum of the first $n$ odd numbers? Linearity gives
\[
\sum_{k=1}^n(2k-1)=2\sum_{k=1}^nk-\sum_{k=1}^n1=n(n+1)-n=n^2.
\]
So $1+3+5+7=16=4^2$. The L-shaped borders in Figure~\ref{fig:gauss} are the same fact drawn for $n=10$: the first $10$ odd numbers fill a square of side $10$. This square is the last step of the third proof of the cube identity.

What is $\sum_{i=1}^n\sum_{j=1}^nij$? The range is a square, not the triangle of Figure~\ref{fig:region}, and the summand is a product $i\cdot j$. The inner sum is $i\sum_{j=1}^nj=i\cdot n(n+1)/2$. The outer sum is then
\[
\frac{n(n+1)}{2}\sum_{i=1}^ni=\left(\frac{n(n+1)}{2}\right)^2.
\]
A product $f(i)g(j)$, summed over a rectangle, factors into a product of two separate sums. The fourth proof of the cube identity starts from this factorisation.
\end{examplex}

\begin{think}
Compute $1+2+4+\dots+2^{10}$ from the geometric formula, and again from the binary reading above. Which argument do you trust more, and why? The formula covers every ratio $q$. The binary reading explains the case $q=2$, and Figure~\ref{fig:geometric} is that case drawn for a shorter bar.
\end{think}

\section{Telescoping}\label{sec:telescope}

Pairing found one closed form because the list could be reversed. A second way to find a closed form is to write the general term as a difference $v_{k+1}-v_k$. The sum then collapses to its two ends, as in Figure~\ref{fig:telescope}. Induction, in the next section, is what you use when you cannot see the cancellation and still want a proof for every $n$.

\begin{figure}[ht]
\centering
\fig{ch01-telescope}
\caption{Each term $1/(k(k+1))$ is a difference $1/k-1/(k+1)$. In the sum up to $k=4$, every intermediate value appears once with a plus and once with a minus. The rust arcs mark those cancellations. Only $1$ and $-1/5$ remain.}
\label{fig:telescope}
\end{figure}

We write $\prod_{k=m}^n a_k$ for the product $a_m\cdots a_n$. An empty product equals $1$, because $1$ is the number that leaves a product unchanged. There was no factor to apply, so the product should stay as it was. The factorial is $n!=1\cdot2\cdots n$ for $n\ge1$, and the same reason sets $0!=1$. Then $(n+1)!=(n+1)\,n!$ also holds at $n=0$, since both sides equal $1$.

\begin{think}
Compute $\frac1{1\cdot2}+\frac1{2\cdot3}+\dots+\frac1{5\cdot6}$ by hand before you read the general formula. The total is a simple fraction. Guess the sum up to $\frac1{n(n+1)}$, and try to write $\frac1{k(k+1)}$ as a difference of two simpler fractions.
\end{think}

\begin{proposition}[telescoping sums]\label{prop:telescope}
If $u_k=v_{k+1}-v_k$ for every $k$ in the range of summation, then $\displaystyle\sum_{k=0}^{n-1}u_k=v_n-v_0$.
\end{proposition}
\begin{proof}
Write the sum out. For $n=0$ there are no terms, and $v_0-v_0=0$. For $n\ge1$,
\[
\sum_{k=0}^{n-1}(v_{k+1}-v_k)=(v_1-v_0)+(v_2-v_1)+\dots+(v_n-v_{n-1}).
\]
Each $v_k$ with $0<k<n$ is added once and subtracted once. Those terms cancel, and the sum is $v_n-v_0$.
\end{proof}

The name comes from a telescope sliding shut. You do not get a closer look at the middle terms. You get rid of them. The same cancellation is the step of a weak induction, one new term at a time. The written sum already shows every term, so the induction can wait until the next section, where the terms are no longer visible in advance.

The work is to find $v$. If $u_k$ is a rational function of $k$, try to split it into fractions whose denominators are shifts of one another. If $u_k$ contains a factorial, try $v_k=k!$ or $v_k=(k+1)!$. If $u_k$ is a product of consecutive integers, lengthen the product: for instance
\[
k(k+1)(k+2)-(k-1)k(k+1)=3k(k+1),
\]
so a product of two consecutive integers is a difference of two products of three, divided by $3$.

\begin{example}
Since $\dfrac{1}{k(k+1)}=\dfrac1k-\dfrac1{k+1}$, Proposition~\ref{prop:telescope} gives
\[
\sum_{k=1}^n\frac1{k(k+1)}=1-\frac1{n+1}.
\]
For $n=4$ this is the $4/5$ of Figure~\ref{fig:telescope}. For $n=0$ both sides are $0$, the empty sum.
\end{example}

\begin{examplex}
First, $\sum_{k=1}^nk\cdot k!=(n+1)!-1$. The difference is $k\cdot k!=(k+1)!-k!$, because $k\cdot k!=(k+1-1)\,k!=(k+1)!-k!$. The sum collapses to $(n+1)!-1!$.

Second, $\sum_{k=1}^nk(k+1)=n(n+1)(n+2)/3$. The difference above rearranges to
\[
k(k+1)=\frac13\bigl(k(k+1)(k+2)-(k-1)k(k+1)\bigr).
\]
Summing from $k=1$ to $n$ leaves only the upper end $n(n+1)(n+2)/3$. Expanding $k(k+1)=k^2+k$ and using the arithmetic-progression formula then gives $\sum k^2$ with no new idea:
\[
\sum_{k=1}^nk^2=\frac{n(n+1)(n+2)}{3}-\frac{n(n+1)}{2}=\frac{n(n+1)(2n+1)}{6}.
\]
Check $n=2$: $1+4=5$, and $2\cdot3\cdot5/6=5$. The binomial basis will give the same formula again, which is a useful check that two tools agree.

Products telescope as well. For $n\ge2$,
\[
\prod_{k=2}^n\left(1-\frac1k\right)=\prod_{k=2}^n\frac{k-1}{k}=\frac{1\cdot2\cdots(n-1)}{2\cdot3\cdots n}=\frac1n.
\]
Every factor cancels except the numerator $1$ and the denominator $n$.
\end{examplex}

Telescoping finds the closed form. It does not only certify a guess you already had. The partial fraction in the first example is the model for Exercise~\ref{ex:tel3}.

\section{Induction, in three flavours}\label{sec:induction}

The progressions and the telescoping sums were proved by writing the terms. Many claims about every natural number do not pair, and do not cancel on the page. We need a principle that starts from $\N$ itself.

Throughout the book, $\N=\{0,1,2,\dots\}$ and $\N^*=\{1,2,\dots\}$. Chapter~\ref{ch:language} recorded the one property of $\N$ we take as an axiom.

\begin{reminder}[well ordering]
Every nonempty subset of $\N$ has a least element.
\end{reminder}

Arguments that prove a first case, and then a step from one case to the next, are older than any single statement of a principle. Al-Karaji, around the year $1000$, and al-Samaw'al in \emph{al-Bahir}, about $1150$, treat the first few exponents of a binomial expansion and say that the same step continues. The table of coefficients they are building is the triangle of the next section, and it survives in al-Samaw'al's account of a lost work of al-Karaji. Levi ben Gershon, in the \emph{Maasei Hoshev} (``The Art of Calculation''), finished in $1321$, uses the same shape for counting claims. Pascal's \emph{Trait\'e du triangle arithm\'etique}, written in $1654$ and printed in $1665$, states a principle that applies to an arbitrary proposition about the natural numbers. What follows is that principle, derived from well ordering, so that a later appeal is an appeal to a proof.

Let $P(n)$ be a statement about $n\in\N$. Well ordering can be rewritten in two other ways. All three say the same thing about $\N$.

\begin{proposition}[induction principles]\label{prop:induction}
\begin{enumerate}[label=(\roman*)]
\item \textbf{Weak induction.} If $P(0)$ holds, and if $P(n)$ implies $P(n+1)$ for every $n\in\N$, then $P(n)$ holds for every $n\in\N$.
\item \textbf{Strong induction.} If, for every $n\in\N$, the truth of $P(k)$ for all $k<n$ implies $P(n)$, then $P(n)$ holds for every $n\in\N$.
\end{enumerate}
Together with well ordering, these are three equivalent principles.
\end{proposition}
\begin{proof}
We derive strong induction from well ordering, weak induction from strong induction, and well ordering from weak induction.

\emph{Well ordering gives strong induction.} Suppose that for every $n$, the statements $P(k)$ with $k<n$ imply $P(n)$. Let $E=\{n\in\N:P(n)\text{ is false}\}$. We claim $E$ is empty. If not, well ordering supplies a least element $n$ of $E$. Every $k<n$ lies outside $E$, so $P(k)$ is true for every $k<n$. The hypothesis then says $P(n)$ is true, so $n\notin E$. That contradicts the choice of $n$. Hence $E$ is empty, and $P(n)$ holds for every $n$.

\emph{Strong induction gives weak induction.} Assume $P(0)$, and assume $P(n)$ implies $P(n+1)$ for every $n$. We check the hypothesis of strong induction. Fix $n$, and assume $P(k)$ for every $k<n$. If $n=0$, there is no natural number $k<0$, so that assumption asks for nothing. The separate hypothesis $P(0)$ gives the conclusion. If $n\ge1$, then $n-1<n$, so $P(n-1)$ holds, and the weak step gives $P(n)$. Strong induction now yields $P(n)$ for every $n$.

\emph{Weak induction gives well ordering.} Let $A\subseteq\N$, and suppose $A$ has no least element. We show $A$ is empty. Let $Q(n)$ say that none of $0,1,\dots,n$ lies in $A$. First, $0\notin A$: otherwise $0$ would be the least element of $A$. So $Q(0)$ holds. If $Q(n)$ holds, then $n+1\notin A$ as well, because $n+1$ would otherwise be the least element of $A$. Thus $Q(n+1)$ holds. Weak induction gives $Q(n)$ for every $n$, so no natural number lies in $A$.
\end{proof}

Read the strong form once more at $n=0$. The phrase ``$P(k)$ for every $k<0$'' has nothing to check, because no such $k$ exists in $\N$. A condition with nothing to check is true. So the inductive step, at $n=0$, must prove $P(0)$ from no earlier case. The base case is still there. It is hidden inside the step.

The well-ordering form is the engine of a proof by \emph{minimal counterexample}: assume some natural number fails, take the least failure, and build a contradiction. The same shape, producing from a failure an even smaller failure, is Fermat's infinite descent. Section~\ref{sec:divisibility} uses the shape on remainders. Here is the shape on factors.

An integer $p>1$ is \emph{prime} when its only positive divisors are $1$ and $p$. A positive integer $d$ \emph{divides} $a$, written $d\mid a$, when $a=dc$ for some integer $c$.

\begin{example}
Every integer $n\ge2$ is a product of primes. Suppose not, and let $n$ be the least integer that is at least $2$ and is not a product of primes. Then $n$ itself is not prime, so $n=ab$ with integers $a,b$ satisfying $1<a<n$ and $1<b<n$. Both $a$ and $b$ are smaller than $n$ and at least $2$, so each is a product of primes. Therefore $n$ is a product of primes. That contradicts the choice of $n$. So no such $n$ exists.
\end{example}

The example uses strong induction in the shape of a least counterexample. Weak induction only remembers the previous integer. Here $12=3\cdot4$, and neither factor is $11$. A proof that can only see $n-1$ cannot see those factors. The uniqueness of the product needs Euclid's lemma, proved from well ordering in Section~\ref{sec:divisibility}. Volume~2, \emph{The Integers: Primes, Patterns and Proofs}, builds the same facts into a theory of primes.

\begin{reminder}[how to write an induction]
A proof by weak induction has three visible parts. Write $P(n)$ in full, including the range of $n$. Check the base case; do not wave at it. In the step, begin from $P(n)$ and end at $P(n+1)$, and use the assumption at least once. If the assumption is never used, the argument did not need induction. If the step only works for $n\ge2$, say so, and check the missing small cases by hand.

The usual false proof that all horses have the same colour fails on exactly this point. Let $P(n)$ say that any set of $n$ horses has a single colour. Then $P(1)$ is true. Given $n+1$ horses, the first $n$ share a colour and the last $n$ share a colour, so all $n+1$ would share a colour, if the two groups of $n$ had a horse in common to pass the colour along. For the step from $1$ horse to $2$, the two groups are disjoint. The step does not start. A later step cannot repair a missing start. Figure~\ref{fig:chain} is that broken first arrow. Often a stronger statement is easier to prove, because it gives a stronger hypothesis to use in the step. Exercise~\ref{ex:sumsq} is the model.
\end{reminder}

\begin{figure}[ht]
\centering
\fig{ch01-chain}
\caption{Weak induction is a chain of steps. In the false proof that all horses share a colour, $P(1)$ holds, and the step from $n$ to $n+1$ needs the two groups of $n$ horses to overlap. From one horse to two, they do not. The rust arrow is the step that never starts. Later arrows cannot repair it.}
\label{fig:chain}
\end{figure}

\begin{examplex}[Golomb's trominoes]
An L-tromino is three squares joined in an L. Remove any one square from a $2^m\times2^m$ board. The rest can be tiled by L-trominoes. Figure~\ref{fig:tromino} shows one stage of the construction.

We induct on $m$. For $m=1$, a $2\times2$ board minus one square is itself one tromino. Suppose every $2^m\times2^m$ board with one square removed can be tiled, and consider a $2^{m+1}\times2^{m+1}$ board with one square removed. Cut the board into four $2^m\times2^m$ quadrants. The missing square lies in one quadrant. Place one tromino at the centre so that it covers one square in each of the other three quadrants. Each quadrant is now a $2^m\times2^m$ board with one square missing. The inductive hypothesis tiles all four.

As a count: the tiling uses $(4^m-1)/3$ trominoes. This is an integer, by a second induction. For $m=1$ the quotient is $1$. The step is the algebra
\[
\frac{4^{m+1}-1}{3}=4\cdot\frac{4^m-1}{3}+1,
\]
so an integer value at $m$ produces an integer value at $m+1$. The number of tiles and the divisibility $3\mid(4^m-1)$ are the same fact.
\end{examplex}

\begin{figure}[ht]
\centering
\fig{ch01-tromino}
\caption{An $8\times8$ board with one square removed, tiled by the inductive step. The dashed lines cut the board into four $4\times4$ quadrants. The central tromino covers one square of each quadrant that does not contain the hole. Each quadrant is then a smaller copy of the same problem. There are $(4^3-1)/3=21$ trominoes.}
\label{fig:tromino}
\end{figure}

The next inequality is the first one the sequences chapter will quote.

\begin{proposition}[Bernoulli's inequality]\label{prop:bernoulli}
For $x\ge-1$ and $n\in\N$, $(1+x)^n\ge1+nx$. The power $n=0$ is $1$, including the case $x=-1$.
\end{proposition}
\begin{proof}
For $n=0$ both sides are $1$. Suppose the inequality holds for some $n$. The number $1+x$ is nonnegative, because $x\ge-1$. Multiplying by a nonnegative number preserves the inequality direction:
\[
(1+x)^{n+1}=(1+x)\,(1+x)^n\ge(1+x)(1+nx)=1+(n+1)x+nx^2.
\]
The extra term $nx^2$ is nonnegative. Dropping it only weakens the right side, so $(1+x)^{n+1}\ge1+(n+1)x$. Weak induction finishes the proof.

The bound $x\ge-1$ is necessary for this argument. If $x<-1$, then $1+x$ is negative, and multiplication would reverse the inequality we were trying to keep.
\end{proof}

Check $n=2$ and $x=1$: $(1+1)^2=4$ and $1+2\cdot1=3$, and $4\ge3$. Check $x=-1$ and $n=3$: the left side is $0$, and the right side is $1-3=-2$.

Jacob Bernoulli printed a comparison of this kind in his \emph{Positiones arithmeticae de seriebus infinitis} (Basel, $1689$). Take a geometric progression and the arithmetic progression that shares its first two terms. His proposition says that each later geometric term is larger than the corresponding arithmetic term. The inequality proved here is the form later arguments quote: it allows $x=-1$, and it allows equality.

\portrait{jbernoulli}{Jacob Bernoulli (1655--1705). Portrait by his brother Nicolaus Bernoulli, 1687. Public domain, via Wikimedia Commons.}

\section{Binomial coefficients}\label{sec:binom}

Some sums count, because each term counts. With binomial coefficients, that count is the whole argument.

\begin{think}
Write the first sixteen rows of Pascal's triangle, and shade the odd entries. What shape appears? Look at Figure~\ref{fig:pascal} only after you have drawn yours. Then try to see why row $2^m-1$ is entirely odd, and why row $2^m$ is odd only at its two ends. Exercise~\ref{ex:pascalmod2} asks for a proof.
\end{think}

The binomial coefficient $\binom nk$ counts the $k$-element subsets of an $n$-element set. It is read ``$n$ choose $k$''. We set $\binom nk=0$ when $k>n$ or $k<0$. A set with $n$ elements has one empty subset and one full subset, so $\binom n0=\binom nn=1$. It has $n$ subsets of size $1$, so $\binom n1=n$. Choosing which $k$ elements to keep is the same decision as choosing which $n-k$ elements to leave out, so $\binom nk=\binom n{n-k}$.

\begin{proposition}\label{prop:binom}
For $n\ge1$ and every integer $k$,
\begin{align*}
\text{(a)}\quad &\binom nk=\binom{n-1}{k-1}+\binom{n-1}{k},\\
\text{(b)}\quad &\binom nk=\frac{n!}{k!\,(n-k)!}\qquad(0\le k\le n),\\
\text{(c)}\quad &(x+y)^n=\sum_{k=0}^n\binom nk x^ky^{n-k}.
\end{align*}
Identity (c) holds for all real $x$ and $y$. The argument uses only associativity, commutativity and distributivity. It therefore applies to the complex numbers of Chapter~\ref{ch:complex}, and to any quantities that commute under multiplication. If the factors do not commute, the order in which they are chosen changes the product, and the terms no longer sort themselves by the number of $x$'s. Volume~3 is where such products become the subject.
\end{proposition}
\begin{proof}
\emph{(a), Pascal's rule.} Let $E$ be a set with $n$ elements, and fix one element $a\in E$. A $k$-element subset either contains $a$ or it does not. If it contains $a$, the rest is a $(k-1)$-element subset of $E\setminus\{a\}$. If it does not, it is a $k$-element subset of $E\setminus\{a\}$. These two kinds do not overlap, and every $k$-element subset is of one kind. That is the two terms on the right. Figure~\ref{fig:split} is the case $n=3$, $k=2$.

\emph{(b), the factorial formula.} Count the ordered sequences of $k$ distinct elements of $E$, in two ways. Choosing them in order gives $n$ choices for the first entry, then $n-1$ for the second, and so on, down to $n-k+1$ choices for the last. The product is
\[
n(n-1)\cdots(n-k+1)=\frac{n!}{(n-k)!}.
\]
Choosing the underlying set first, in $\binom nk$ ways, and then ordering it, in $k!$ ways, gives $k!\binom nk$. The two counts are equal. For $k=0$ both sides of (b) are $1$, since $0!=1$ and there is one empty sequence.

\emph{(c), the binomial theorem.} Expand the product of $n$ factors $(x+y)\cdots(x+y)$. Each term in the expansion comes from choosing either $x$ or $y$ in every factor, then multiplying the choices. The terms that choose $x$ in exactly $k$ factors, and $y$ in the other $n-k$ factors, all equal $x^ky^{n-k}$, once the factors are allowed to change order. The number of ways to choose which $k$ factors contribute $x$ is $\binom nk$.
\end{proof}

\begin{figure}[ht]
\centering
\fig{ch01-split}
\caption{Pascal's rule for $\binom{3}{2}$. The charcoal box holds the subsets that contain a fixed element $a$. The rust box holds the subsets that omit $a$. The two boxes are disjoint and exhaust the subsets of size $2$, so $3=2+1$.}
\label{fig:split}
\end{figure}

Pascal's rule is the recipe for the triangle: each entry is the sum of the two entries above it. Formula (b) is the one to compute with. For instance
\[
\binom{10}{3}=\frac{10\cdot9\cdot8}{3\cdot2\cdot1}=120,
\]
where the factorial $7!$ was cancelled before any multiplication. Cancel before you multiply. The fraction is a count of subsets, and a count is an integer. If a prime in the denominator has not cancelled, look again: either the range is wrong, or a factor is still hiding in the numerator.

Formula (c) is why these numbers are called binomial coefficients. Two cases are worth remembering. With $x=y=1$,
\[
\sum_{k=0}^n\binom nk=2^n,
\]
because a set with $n$ elements has $2^n$ subsets: each element is in or out. With $x=-1$ and $y=1$, and with $n\ge1$,
\[
\sum_{k=0}^n(-1)^k\binom nk=0.
\]
A set with at least one element has as many subsets of even size as of odd size. For $n=0$ the same sum is $1$, not $0$.

The triangle is older than Pascal's treatise, and it was used for different jobs. Yang Hui's \emph{Xiangjie jiuzhang suanfa} ($1261$) prints the array and attributes the method to Jia Xian, in the eleventh century, who used the entries to extract roots. Figure~\ref{fig:rows} is the first rows of that array. Al-Samaw'al's \emph{al-Bahir}, about $1150$, preserves a table through the twelfth power and attributes the construction to al-Karaji. Omar Khayyam, later in the eleventh century, also used these coefficients to extract roots (Edwards, in the sources). Zhu Shijie's \emph{Jade Mirror of the Four Unknowns}, printed in $1303$, prints the triangle again. Pascal's treatise organises the consequences, including an explicit induction principle of the kind derived in the previous section, and it is the printed page in the portrait below. The counting proof above does not use any of those dates.

\begin{figure}[ht]
\centering
\fig{ch01-rows}
\caption{The first rows of the binomial triangle, drawn for this chapter. The circled entry $10$ is the sum of the two rust entries above it, $4$ and $6$. That is Pascal's rule, read on the array. Yang Hui's book of $1261$ prints such an array and attributes the method to Jia Xian.}
\label{fig:rows}
\end{figure}

\portraits{pascal}{Blaise Pascal (1623--1662). Anonymous copy of a portrait by Fran\c cois Quesnel, c.\,1690, Palace of Versailles. Public domain, via Wikimedia Commons.}{pascal-triangle}{Pascal's arithmetical triangle, as printed in his \emph{Trait\'e du triangle arithm\'etique} (written 1654, published 1665). Public domain, via Wikimedia Commons.}

\begin{figure}[ht]
\centering
\fig{ch01-pascal}
\caption{Left: Pascal's triangle modulo $2$. The odd entries, drawn in charcoal, form a Sierpi\'nski triangle. Exercise~\ref{ex:pascalmod2} explains the rows of length a power of two. Lucas' theorem, an exercise in the congruences chapter of Volume~2, explains every row. Right: the $5\times5$ multiplication table, cut into L-shaped regions $L_k=\{(i,j):\max(i,j)=k\}$. Each region sums to a cube, and the whole table sums to $(1+\dots+5)^2$. The fourth proof in Section~\ref{sec:fourproofs} writes the picture out.}
\label{fig:pascal}
\end{figure}

Most binomial identities in this chapter are proved in one of two ways. \emph{Double counting} counts one set in two ways, as in (b). \emph{Comparing coefficients} reads one polynomial in two writings, as in (c). The next identity is both.

\begin{think}
In how many ways can you choose a committee of $2$ people from $3$ women and $2$ men? Count the committees directly. Then count them by the number of women on the committee. Which identity have you proved for these small numbers?
\end{think}

\begin{proposition}[Vandermonde]\label{prop:vandermonde}
For $m,n,r\in\N$,
\[
\sum_{k=0}^{r}\binom mk\binom n{r-k}=\binom{m+n}{r}.
\]
\end{proposition}
\begin{proof}
\emph{By counting.} Choose a committee of $r$ people from $m$ women and $n$ men. There are $\binom{m+n}{r}$ such committees. Group them by the number $k$ of women. A committee with exactly $k$ women is chosen by picking those $k$ women and the remaining $r-k$ men, in $\binom mk\binom n{r-k}$ ways. Sum over $k$. Terms with $k>m$ or $r-k>n$ are zero by the convention on binomial coefficients, so the limits $0$ and $r$ do not need trimming.

\emph{By coefficients.} Compare the coefficient of $x^r$ in $(1+x)^m(1+x)^n=(1+x)^{m+n}$. The right side contributes $\binom{m+n}{r}$. The left side contributes $\sum_k\binom mk\binom n{r-k}$.
\end{proof}

The identity is often called Vandermonde's convolution, after Alexandre-Th\'eophile Vandermonde's paper of $1772$. The same convolution appears in Zhu Shijie's \emph{Jade Mirror of the Four Unknowns} ($1303$). The committee argument above does not depend on either date. It depends on the two counts being counts of the same set. Multiplying $(1+x)^m(1+x)^n$ and reading the coefficient of $x^r$ is the same identity, and it is Exercise~\ref{ex:squares}. The next chapter counts this way again.

\section{Polynomials in the binomial basis}\label{sec:basis}

Binomial coefficients began as counts. The same formulas, read as polynomials in the upper index, are coordinates for every polynomial. That is why a sum of polynomial values has a closed form.

If every polynomial in $k$ is a combination of $\binom k0,\binom k1,\binom k2,\dots$, then every polynomial can be summed by summing those binomial polynomials. The sum of $\binom kr$ along $k$ is the hockey-stick identity, drawn in Figure~\ref{fig:hockey} and proved in the corollary below. The recipe for the coordinates is Newton's. He used finite differences to interpolate, in Book~III, Lemma~V of the \emph{Principia} ($1687$), in the part of that book that treats the motion of comets.

\begin{figure}[ht]
\centering
\fig{ch01-hockey}
\caption{The hockey-stick identity, drawn in Pascal's triangle. The charcoal entries $\binom k2$ for $k=2,\dots,6$ are the handle. They add to the rust entry $\binom73$, the blade, just beyond the last term. Pascal's rule, read upwards, is the reason. Corollary~\ref{cor:sumpoly} proves the general sum, and Exercise~\ref{ex:hockey} asks for a counting proof.}
\label{fig:hockey}
\end{figure}

A \emph{polynomial} is a finite sum $a_0+a_1X+\cdots+a_dX^d$. The numbers $a_j$ are its coefficients. Its degree is the largest index with a nonzero coefficient. Two polynomials are equal when corresponding coefficients agree. The zero polynomial has no largest nonzero index, so its degree is left undefined here. We work with real coefficients. The same arguments apply to rational coefficients, and later to complex coefficients. For a polynomial $P$, the \emph{forward difference} is $\Delta P(X)=P(X+1)-P(X)$, and $\Delta^kP=\Delta(\Delta^{k-1}P)$, with $\Delta^0P=P$.

If $\deg P=d\ge1$ and the leading term is $aX^d$, then the leading terms of $P(X+1)$ and $P(X)$ agree, so they cancel in the difference. Hence $\deg\Delta P<d$. The difference of a constant is $0$.

The binomial polynomials are
\[
\binom Xk=\frac{X(X-1)\cdots(X-k+1)}{k!}\qquad(k\ge0),
\]
with $\binom X0=1$. The product in the numerator has $k$ factors, so the degree is $k$. At each integer $n\ge0$ the value is the binomial coefficient $\binom nk$. At $n=0,1,\dots,k-1$ the value is $0$, because one factor in the numerator vanishes.

\begin{proposition}\label{prop:deltabinom}
For $k\ge1$, $\Delta\binom Xk=\binom X{k-1}$. Also $\Delta\binom X0=0$.
\end{proposition}
\begin{proof}
For $k=1$, $(X+1)-X=1=\binom X0$. For $k\ge2$, factor the product that the two binomial polynomials share:
\begin{align*}
\binom{X+1}k-\binom Xk
&=\frac{X(X-1)\cdots(X-k+2)}{k!}\bigl((X+1)-(X-k+1)\bigr)\\
&=\frac{X(X-1)\cdots(X-k+2)}{k!}\cdot k\\
&=\frac{X(X-1)\cdots(X-k+2)}{(k-1)!}
=\binom X{k-1}.
\end{align*}
The middle step is the arithmetic $(X+1)-(X-k+1)=k$. Finally $\binom X0=1$, whose difference is $0$. This is Pascal's identity, read as an identity of polynomials. It does not need a theorem about roots.
\end{proof}

\begin{theorem}[Newton's forward difference formula]\label{thm:newton}
Every nonzero polynomial $P$ of degree $d$ can be written in exactly one way as
\[
P(X)=\sum_{k=0}^dc_k\binom Xk,
\]
and the coefficients are $c_k=(\Delta^kP)(0)$.
\end{theorem}
\begin{proof}
We first show that the coefficients exist and are unique. We then identify them.

\emph{Existence and uniqueness.} The polynomials $\binom X0,\dots,\binom Xd$ have degrees $0,1,\dots,d$. Only $\binom Xd$ has degree $d$. Match the leading coefficient of $P$ against the leading coefficient of $c_d\binom Xd$. That fixes $c_d$. Subtract $c_d\binom Xd$. The difference has degree less than $d$, or it is zero. Repeat. You reach a constant, which is $c_0\binom X0$. If two such writings existed, their difference would be a writing of the zero polynomial. Comparing leading coefficients from the top degree downward forces every coefficient in that difference to be $0$. So the writing is unique. In the language of Chapter~\ref{ch:linear}, these polynomials form a basis of the polynomials of degree at most $d$. The argument here does not need that chapter. It only uses degrees.

\emph{The formula for $c_k$.} Apply $\Delta^j$ to both sides. Proposition~\ref{prop:deltabinom}, used $j$ times, moves the lower index down by $j$:
\[
\Delta^jP=\sum_{k\ge j}c_k\binom X{k-j}.
\]
Set $X=0$. If $k>j$, then $\binom{0}{k-j}=0$. The remaining term is $c_j\binom00=c_j$. So $(\Delta^jP)(0)=c_j$.
\end{proof}

You do not need to expand $\Delta^kP$ as a polynomial. Write the values $P(0),P(1),P(2),\dots$ in a row. Write the differences of consecutive entries in the row below, and repeat. The first entry of row $k$ is $(\Delta^kP)(0)$. Figure~\ref{fig:differences} does this for $P(X)=X^3$. The first entries are $0,1,6,6$, and the next row is $0$, so
\[
X^3=\binom X1+6\binom X2+6\binom X3.
\]
You can check the identity by expanding the right side. Section~\ref{sec:fourproofs} sums it.

\begin{figure}[ht]
\centering
\fig{ch01-differences}
\caption{The difference table of $P(k)=k^3$. Each lower row holds the differences of consecutive entries above it. The boxed first entries are the coefficients of $k^3$ in the binomial basis. The third differences are constant, and the fourth differences are zero, as they must be for a polynomial of degree $3$.}
\label{fig:differences}
\end{figure}

\begin{corollary}[sums of polynomials]\label{cor:sumpoly}
If $P=\sum_{k=0}^dc_k\binom Xk$ with $c_d\ne0$, then for every $n\ge0$
\[
\sum_{j=0}^{n-1}P(j)=\sum_{k=0}^dc_k\binom n{k+1}.
\]
In particular, the sum is a polynomial in $n$ of degree $d+1$, with leading coefficient $c_d/(d+1)!$.
\end{corollary}
\begin{proof}
The engine is the identity $\sum_{j=0}^{n-1}\binom jk=\binom n{k+1}$. This is the hockey stick of Figure~\ref{fig:hockey}, in a range that starts at $0$. Exercise~\ref{ex:hockey} asks for a second proof, by counting. Here is the inductive proof. For $n=0$ both sides are $0$. Assume the identity for some $n$. The sum up to $n$, rather than $n-1$, adds the single term $\binom nk$. Pascal's rule says
\[
\binom{n+1}{k+1}-\binom n{k+1}=\binom nk,
\]
so the right side grows by exactly that term. This is the identity at $n+1$.

Weight the identity by $c_k$ and sum over $k$. That is the formula for $\sum P(j)$. The polynomial $\binom n{k+1}$ has degree $k+1$ and leading coefficient $1/(k+1)!$, because the numerator is a product of $k+1$ terms $n(n-1)\cdots(n-k)$ and the denominator is $(k+1)!$. The highest degree in the sum comes from $k=d$, and $c_d\ne0$, so the leading coefficient is $c_d/(d+1)!$.
\end{proof}

\begin{examplex}
From the table in Figure~\ref{fig:differences},
\[
\sum_{j=0}^{n-1}j^3=\binom n2+6\binom n3+6\binom n4.
\]
Replacing $n$ by $n+1$ is the second proof in Section~\ref{sec:fourproofs}. For squares, the difference table of $k^2$ begins $0,1,2,0$, so $k^2=\binom k1+2\binom k2$ and
\[
\sum_{j=0}^{n}j^2=\binom{n+1}{2}+2\binom{n+1}{3}=\frac{n(n+1)(2n+1)}{6}.
\]
This agrees with the telescoping computation above. The same table sums $k^4$, $k^5$, or $k^3-2k+7$ in a few lines. Johann Faulhaber, in $1631$, carried the method by hand as far as the seventeenth power. The challenges ask for the fifth, and for why the odd powers are special.
\end{examplex}

Once the difference table is written, Corollary~\ref{cor:sumpoly} sums the polynomial. You no longer have to guess the closed form. The cubes are the first sum we do that way.

\section{Four proofs of one identity}\label{sec:fourproofs}

Now the sentence from the first page, with each of those tools.

\begin{propositionx}[sum of the first cubes]
For every integer $n\ge0$,
\[
\sum_{k=1}^n k^3=\left(\frac{n(n+1)}{2}\right)^2.
\]
For $n=0$ both sides are $0$.
\end{propositionx}

The right side is the square of the $n$th triangular number. The four proofs below are four explanations of that square.

\begin{proof}[First proof: induction]
The formula at $n+1$ differs from the formula at $n$ by the single cube $(n+1)^3$. Add that cube to the inductive square and factor. The next square appears.

For $n=0$ both sides are $0$. For $n=1$ both sides equal $1$. Assume the formula for some $n\ge0$. Then
\begin{align*}
\sum_{k=1}^{n+1}k^3
&=\left(\frac{n(n+1)}{2}\right)^2+(n+1)^3\\
&=\frac{n^2(n+1)^2}{4}+(n+1)^3\\
&=\frac{(n+1)^2}{4}\bigl(n^2+4(n+1)\bigr)\\
&=\frac{(n+1)^2}{4}\bigl(n^2+4n+4\bigr)\\
&=\frac{(n+1)^2(n+2)^2}{4}.
\end{align*}
The last step is $n^2+4n+4=(n+2)^2$. So the formula holds for $n+1$. The hypothesis was used in the first line, and the algebra after that is factoring. For a check beyond the base, $n=4$ gives $100$ on both sides, as in the opening.

The step starts at $n=0$, and it uses the hypothesis. That is the repair the horse-colour argument never made. When a sum grows by a term you can already see in the closed form, this is the proof that fits.
\end{proof}

\begin{proof}[Second proof: a change of basis]
Figure~\ref{fig:differences} gave $k^3=\binom k1+6\binom k2+6\binom k3$. Expand the right side once, so the identity is not only a reading of a table:
\begin{align*}
\binom k1+6\binom k2+6\binom k3
&=k+3k(k-1)+k(k-1)(k-2)\\
&=k+3k^2-3k+k^3-3k^2+2k\\
&=k^3.
\end{align*}
The first line uses $6\binom k2=3k(k-1)$ and $6\binom k3=k(k-1)(k-2)$. The second line expands $k(k-1)(k-2)=k(k^2-3k+2)=k^3-3k^2+2k$. The cubic terms are $k^3$, the quadratic terms $3k^2-3k^2$ cancel, and the linear terms $k-3k+2k$ cancel. Sum from $k=1$ to $n$. Corollary~\ref{cor:sumpoly}, applied with upper index $n+1$, says
\[
\sum_{k=0}^{n}\binom kr=\binom{n+1}{r+1}.
\]
For $r\ge1$ the $k=0$ term is $0$, so the sum may start at $k=1$. Therefore
\begin{align*}
\sum_{k=1}^n k^3
&=6\binom{n+1}{4}+6\binom{n+1}{3}+\binom{n+1}{2}.
\end{align*}
Factor $(n+1)n/4$ out of each term. The three binomial coefficients contribute
\[
(n-1)(n-2),\qquad 4(n-1),\qquad 2,
\]
because
\begin{align*}
6\binom{n+1}{4}&=\frac{(n+1)n(n-1)(n-2)}{4},\\
6\binom{n+1}{3}&=(n+1)n(n-1)=\frac{(n+1)n}{4}\cdot4(n-1),\\
\binom{n+1}{2}&=\frac{(n+1)n}{2}=\frac{(n+1)n}{4}\cdot2.
\end{align*}
Adding the three contributions gives
\[
(n-1)(n-2)+4(n-1)+2=n^2-3n+2+4n-4+2=n^2+n=n(n+1).
\]
Hence the sum equals
\[
\frac{(n+1)n}{4}\cdot n(n+1)=\left(\frac{n(n+1)}{2}\right)^2.
\]
For a check at $n=3$, the binomial form is $\binom42+6\binom43+6\binom44=6+24+6=36$, and $1^3+2^3+3^3=36$.

The same table sums fifth powers. The challenges ask why the odd powers are special. That question is Faulhaber's.
\end{proof}

\begin{proof}[Third proof: odd borders]
Nicomachus, in Book~II, Chapter~20 of the \emph{Introduction to Arithmetic}, grouped the odd numbers as $1$, then $3+5$, then $7+9+11$, then $13+15+17+19$. Those groups add to $1$, $8$, $27$, $64$, the first cubes. He stated the grouping. The identity for every $n$ is the argument below. The right half of Figure~\ref{fig:gauss} draws the same borders as areas.

We claim that for $k\ge1$, the $k$ consecutive odd numbers
\[
k^2-k+1,\quad k^2-k+3,\quad \dots,\quad k^2+k-1
\]
add to $k^3$. They differ by $2$. There are $k$ of them, because the last minus the first is $(k^2+k-1)-(k^2-k+1)=2k-2$, and $(2k-2)/2+1=k$. Each is odd. The product $k(k-1)$ is even, so $k^2-k=2m$ for an integer $m$, and $k^2-k+1=2m+1$. That is the definition of odd. This is an arithmetic progression of $k$ terms, with first term $k^2-k+1$ and difference $2$. The formula of Section~\ref{sec:sums} gives
\[
k(k^2-k+1)+2\cdot\frac{k(k-1)}{2}=k(k^2-k+1)+k(k-1)=k\bigl(k^2-k+1+k-1\bigr)=k^3.
\]

The last number of block $k$ is $k^2+k-1$. The first number of block $k+1$ is $(k+1)^2-(k+1)+1=k^2+k+1$, the next odd number. So the blocks for $k=1,\dots,n$, placed end to end, are exactly the first $N=n(n+1)/2$ odd numbers. The first example of this chapter says that those odd numbers add to $N^2$. Therefore $\sum_{k=1}^n k^3=N^2$.

For $k=4$ the block is $13+15+17+19=64=4^3$, and the four blocks together are the first $10$ odd numbers, which add to $10^2$. That is the square in Figure~\ref{fig:gauss}.

The square is no longer a surprise. The cubes are consecutive odd numbers in disguise, and consecutive odd numbers fill a square.
\end{proof}

\begin{proof}[Fourth proof: the multiplication table]
The $n\times n$ multiplication table has entries $ij$. Their total is $\bigl(\sum_{i=1}^n i\bigr)^2$, by the rectangular factorisation in Section~\ref{sec:sums}. Cut the table into the regions $L_k=\{(i,j):\max(i,j)=k\}$, drawn for $n=5$ on the right of Figure~\ref{fig:pascal}.

The region $L_k$ consists of the cells $(k,j)$ for $1\le j\le k$, and the cells $(i,k)$ for $1\le i\le k-1$. The corner $(k,k)$ is in the first list only, so it is not counted twice. The sum of the entries is
\begin{align*}
\sum_{j=1}^k k\cdot j+\sum_{i=1}^{k-1}i\cdot k
&=k\cdot\frac{k(k+1)}{2}+k\cdot\frac{(k-1)k}{2}\\
&=k\cdot\frac{k^2+k+k^2-k}{2}\\
&=k^3.
\end{align*}
The regions $L_1,\dots,L_n$ partition the table. Adding the regions therefore gives $\sum_{k=1}^n k^3=\bigl(\sum_{k=1}^n k\bigr)^2$.

For $k=2$ the cells are $2\cdot1$, $2\cdot2$ and $1\cdot2$, and $2+4+2=8=2^3$. The corner $4$ was not counted twice.

The square is the table of products, and each cube is one L-shaped piece of it. Same area, two orders of adding.
\end{proof}

Induction, a change of basis, a grouping, and a double count: four proofs of one square. A new closed form usually yields to one of them. A sum that resists all four may be one of the sums that Volumes~6 and~14 study without a closed form.

\section{Divisibility from well ordering}\label{sec:divisibility}

The exercises on Fermat, Fibonacci numbers and harmonic sums use a few facts about divisibility. We prove them here, from well ordering, so that later appeals are appeals to a proof. The existence of a prime factorisation was the minimal-counterexample example in Section~\ref{sec:induction}. This section adds the remainder, B\'ezout's identity, and the uniqueness. Volume~2 organises the same facts. The arguments below are the ones the present chapter uses.

For integers $a$ and $b$, write $a\mid b$ when $b=ac$ for some integer $c$. The integers $a$ and $b$ are \emph{coprime} when $\gcd(a,b)=1$, where $\gcd(a,b)$ denotes their greatest positive common divisor. Set $\gcd(a,0)=|a|$ for $a\ne0$. The symbol $\gcd(0,0)$ is left undefined: every integer divides $0$, so there is no greatest positive common divisor. For $m\ge1$, the congruence $a\equiv b\pmod m$ means $m\mid(a-b)$. These are definitions. The cancellation rules need proofs.

\begin{propositionx}[division with remainder]
Let $a$ be an integer and let $b$ be a positive integer. There are unique integers $q$ and $r$ such that $a=qb+r$ and $0\le r<b$.
\end{propositionx}
\begin{proof}
We find one such pair, then show there is only one.

\emph{Existence.} Consider the nonnegative integers of the form $a-qb$, as $q$ runs through $\Z$. The set is nonempty. The choice $q=-|a|$ gives $a+|a|b$. Since $b\ge1$, we have $|a|b\ge|a|$, so $a+|a|b\ge a+|a|\ge0$. Let $r$ be the least element of the set, and write $r=a-qb$. Then $r\ge0$. If $r\ge b$, the integer $r-b=a-(q+1)b$ would be nonnegative and strictly smaller than $r$, contradicting the choice of $r$. So $0\le r<b$.

\emph{Uniqueness.} Suppose $a=qb+r=q'b+r'$ with $0\le r,r'<b$. Then $r-r'=(q'-q)b$, so $b$ divides $r-r'$. But $|r-r'|<b$. The only multiple of $b$ strictly between $-b$ and $b$ is $0$. Hence $r=r'$ and $q=q'$.

We write this quotient as $\lfloor a/b\rfloor$. The floor of a general real number waits for the Archimedean property, in Chapter~\ref{ch:sequences}. Figure~\ref{fig:division} is the picture: the multiples of $b$ cut the line into steps of length $b$, and $a$ falls in one step.
\end{proof}

\begin{figure}[ht]
\centering
\fig{ch01-division}
\caption{Division of a positive $a$ by $b>0$. The remainder is the distance from $a$ back to the last multiple of $b$ that does not pass $a$, so $0\le r<b$. The next multiple is too far to be the left endpoint.}
\label{fig:division}
\end{figure}

For a check, divide $-7$ by $3$. The proposition promises a nonnegative remainder. Take $q=-3$ and $r=2$: then $-7=(-3)\cdot3+2$, and $0\le2<3$. The pair $q=-2$, $r=-1$ uses the same $a$ and $b$, but the remainder is negative, so it is not the pair in the statement.

\begin{propositionx}[B\'ezout's identity]
Let $a$ and $b$ be integers, not both zero. The least positive integer of the form $au+bv$, with $u,v\in\Z$, equals $\gcd(a,b)$. In particular $\gcd(a,b)=au+bv$ for some integers $u$ and $v$, and every common divisor of $a$ and $b$ divides $\gcd(a,b)$.
\end{propositionx}
\begin{proof}
The positive integers of the form $au+bv$ form a nonempty set: if $a\ne0$, then $|a|$ is one of them, and otherwise $|b|$ is one of them. Let $d$ be the least such positive integer, and write $d=au+bv$.

Divide $a$ by $d$: $a=qd+r$ with $0\le r<d$. Then $r=a-qd=a(1-qu)+b(-qv)$, so $r$ is an integer combination of $a$ and $b$. If $r>0$, then $r$ would be a smaller positive combination than $d$. So $r=0$, and $d\mid a$. The same argument gives $d\mid b$.

Thus $d$ is a positive common divisor. Every common divisor of $a$ and $b$ divides $au+bv$, hence divides $d$. So no common divisor is larger than $d$, and $d=\gcd(a,b)$.
\end{proof}

Here is the consequence the polynomials chapter will quote. Let $p$ be prime, and suppose $p$ does not divide $a$. Then $\gcd(a,p)=1$, so $au+pv=1$ for some integers $u$ and $v$. Reduce modulo $p$: $au\equiv1\pmod p$. Every residue not divisible by $p$ therefore has a multiplicative inverse modulo $p$.

If $a$ and $b$ are positive, their least common multiple exists for the same reason. Their product is a positive common multiple, so the nonempty set of positive common multiples has a least element, by well ordering.

\begin{propositionx}[Euclid's lemma and unique factorisation]
\begin{enumerate}[label=(\alph*)]
\item If $\gcd(a,b)=1$ and $a\mid bc$, then $a\mid c$.
\item If a prime $p$ divides a product of integers, then $p$ divides one of the factors.
\item Every integer $n\ge2$ can be written as a product of primes. The primes in any two such writings are the same, up to order.
\end{enumerate}
\end{propositionx}
\begin{proof}
\emph{Coprime cancellation.} By B\'ezout's identity, $au+bv=1$ for some integers $u$ and $v$. Multiply by $c$: $auc+bvc=c$. The hypothesis $a\mid bc$ says $a$ divides $bvc$. Also $a$ divides $auc$. So $a$ divides the sum, which is $c$.

\emph{A prime divides a factor.} Let $p$ be prime and suppose $p\mid bc$. If $p\mid b$, we are done. If not, then $\gcd(p,b)=1$, because the only positive divisors of $p$ are $1$ and $p$. Part (a) gives $p\mid c$. For a longer product, repeat the same step, one factor at a time.

\emph{Existence of a prime factorisation.} This is the minimal-counterexample argument in Section~\ref{sec:induction}.

\emph{Uniqueness.} Suppose $n\ge2$ has two prime factorisations. A prime $p$ from the first divides the product on the second side, so by (b) it equals one prime on the second side. Cancel that prime. The two shorter products are equal. Repeat. The lists cannot have different lengths: if one list were used up first, a product of one or more primes would equal $1$, which is impossible. So the lists agree up to order.
\end{proof}

In particular, the least common multiple of positive integers is found by taking, for each prime, the largest exponent with which that prime appears in any of them. These facts justify the cancellation and parity arguments in the exercises.

\section{Workshop: a pattern is a suspect, not a verdict}

The sequence $1,2,4,8,16$ looks like the powers of two. It is also the start of many other rules. A finite list cannot say which rule comes next. Five correct values are an alibi. They are not a verdict.

\begin{examplex}[five agreements, then a disagreement]
Let $Q(X)=\sum_{j=0}^4\binom Xj$, with the binomial polynomials of this chapter. For an integer $n$ with $0\le n\le4$, every term with $j>n$ vanishes, and the binomial theorem gives $Q(n)=2^n$. But
\[
Q(5)=1+5+10+10+5=31,
\]
which is not $32$. For any constant $c$, the polynomial
\[
R_c(X)=Q(X)+c\,X(X-1)(X-2)(X-3)(X-4)
\]
takes the same five values at $0,1,2,3,4$, and $R_c(5)=31+120c$, because $5\cdot4\cdot3\cdot2\cdot1=120$. Every proposed sixth value is possible. A pattern can suggest a theorem. Induction, a bijection, or an algebraic identity has to finish the proof. Figure~\ref{fig:pattern} draws the first disagreement.
\end{examplex}

\begin{figure}[ht]\centering\fig{review01-pattern}
\caption{Five exact agreements do not determine the next term. Squares show $2^n$. Circles show $Q(n)=\sum_{j=0}^4\binom nj$. The values agree for $n=0,1,2,3,4$, and then $Q(5)=31$ while $2^5=32$. The lines only guide the eye between the integer inputs.}
\label{fig:pattern}
\end{figure}

\subsection*{Try it, then take a hint}
\textbf{Mission.} Find $\sum_{k=1}^n k(k+1)$ without expanding the summand into $k^2+k$. Then explain why the answer is an integer.

\textbf{Hint 1.} A product of consecutive integers wants a factorial in the denominator. \textbf{Hint 2.} Note that $k(k+1)=2\binom{k+1}{2}$. Pascal's identity then turns the sum into a telescope.

\textbf{Solution.} Proposition~\ref{prop:deltabinom} says $\binom{k+2}{3}-\binom{k+1}{3}=\binom{k+1}{2}$. Therefore
\[
\sum_{k=1}^n k(k+1)=2\sum_{k=1}^n\binom{k+1}{2}=2\binom{n+2}{3}=\frac{n(n+1)(n+2)}{3}.
\]
The binomial coefficient is an integer, so the sum is an integer. Alternatively, among three consecutive integers $n$, $n+1$ and $n+2$, one is divisible by $3$. For $n=0$ the empty sum and the formula both give $0$. For $n=2$ both sides equal $2+6=8$, and $2\cdot3\cdot4/3=8$.

\begin{reminder}[a proof-writing checklist]
Name the domain of the index. Check the first admissible value. State what the induction hypothesis allows, and what it does not. Take the step for an arbitrary index in that domain. Check that every denominator is nonzero. A hundred correct examples establish a hundred examples. One correct inductive step, with its base, establishes the rest.
\end{reminder}

\section*{What this chapter prepares}
The sentence ``for every positive integer $n$'' now has a method behind it. Pairing gave the sum of the first $n$ integers. Induction, the binomial basis, Nicomachus's borders and the multiplication table are the four proofs of the cube identity in the epigraph. Telescoping finds a closed form when the general term is a difference. The binomial basis does the same for every polynomial.

The binomial coefficient has a second job, and the next chapter takes it up. It counts. The question changes from ``what is the sum?'' to ``how many objects are there?'', especially when the restrictions overlap. In how many ways can every guest receive the wrong hat? Pierre R\'emond de Montmort asked the same question of a shuffled pack, and the next chapter counts the arrangements before asking what the count says about chance. Adding one restriction at a time counts some arrangements more than once. The correction is inclusion--exclusion. The same binomial numbers will count lattice paths, by choosing when the path steps east. The same induction will build recurrences. The count comes first. What the count says about chance waits for Chapter~\ref{ch:probability}.

The divisibility above is enough for the exercises. Volume~2, \emph{The Integers: Primes, Patterns and Proofs}, asks what those facts say about the primes. Later, Volume~5 returns to binomial series as functions, with a remainder attached, and Volumes~6 and~14 meet sums that no longer have closed forms.

\section{Exercises}
The identities below use only this chapter, except where a later chapter is named for a second reading. Exercises~\ref{ex:cubes}, \ref{ex:hockey}, \ref{ex:odds} and~\ref{ex:cubesgeo} ask you to write the four proofs of Section~\ref{sec:fourproofs} in your own words, and then to compare the writing with that section. The limits in Exercises~\ref{ex:tel3} and~\ref{ex:telprod} are first examples to revisit after the definitions in Chapter~\ref{ch:sequences}. The finite formulas already show the size of the error.

\begin{exercise}[$\star$]\label{ex:cubes}
For $n\ge0$, prove by induction that $\displaystyle\sum_{k=1}^n k^3=\left(\frac{n(n+1)}{2}\right)^2$.
\end{exercise}

\begin{exercise}[$\star$]\label{ex:tel3}
For $n\ge0$, compute $\displaystyle S_n=\sum_{k=1}^n\frac{1}{k(k+1)(k+2)}$, and then its limit as $n\to\infty$.
\end{exercise}

\begin{exercise}[$\star$]\label{ex:kbinom}
For $n\ge1$, show that $\displaystyle\sum_{k=0}^n k\binom nk=n\,2^{n-1}$ in two ways: by counting, and by the identity $k\binom nk=n\binom{n-1}{k-1}$. Check separately that the sum is $0$ for $n=0$. After the calculus chapter, recover the identity by differentiating $(1+x)^n$.
\end{exercise}

\begin{exercise}[$\star\star$]\label{ex:hockey}
\emph{Hockey stick.} Show that $\displaystyle\sum_{k=r}^n\binom kr=\binom{n+1}{r+1}$ for $0\le r\le n$. Deduce a second proof of the formula for $\sum_{k=1}^n k^3$, using $k^3=6\binom k3+6\binom k2+\binom k1$.
\end{exercise}

\begin{exercise}[$\star\star$]\label{ex:squares}
Show that $\displaystyle\sum_{k=0}^n\binom nk^2=\binom{2n}n$. Deduce that $\binom{2n}n\ge 4^n/(n+1)$.
\end{exercise}

\begin{exercise}[$\star\star$]\label{ex:fermat}
Let $p$ be prime.
\begin{enumerate}[label=(\alph*)]
\item Show that $p$ divides $\binom pk$ for $0<k<p$.
\item Deduce Fermat's little theorem in the form $a^p\equiv a\pmod p$, for every $a\in\N$.
\end{enumerate}
\end{exercise}

\begin{exercise}[$\star\star$]\label{ex:fibgcd}
Let $(F_n)$ be the Fibonacci sequence, with $F_0=0$, $F_1=1$, and $F_{n+2}=F_{n+1}+F_n$.
\begin{enumerate}[label=(\alph*)]
\item Show that $F_{m+n}=F_m F_{n+1}+F_{m-1}F_n$ for $m\ge1$ and $n\ge0$.
\item Show that $\gcd(F_n,F_{n+1})=1$ for $n\ge0$.
\item Show that $\gcd(F_m,F_n)=F_{\gcd(m,n)}$ for $m,n\ge0$ not both zero.
\end{enumerate}
\end{exercise}

\begin{exercise}[$\star\star$]\label{ex:cubesgeo}
\emph{A picture, written out.} The $n\times n$ multiplication table has entries $ij$. Their total is $\bigl(\sum_{i=1}^n i\bigr)^2$. Cut the table into the regions $L_k=\{(i,j):\max(i,j)=k\}$. Show that the entries of $L_k$ sum to $k^3$, and conclude that $\sum_{k=1}^n k^3=\bigl(\sum_{k=1}^n k\bigr)^2$.
\end{exercise}

\begin{exercise}[$\star\star\star$]\label{ex:harmonic}
Let $H_n=\sum_{k=1}^n 1/k$. Show that $H_n$ is not an integer for any $n\ge2$. \emph{Hint: look at the largest power of $2$ that does not exceed $n$.}
\end{exercise}

\subsection*{Further exercises}

\begin{exercise}[$\star$]\label{ex:sumsq}
\emph{Choosing the statement.} Try to prove by induction that $\sum_{k=1}^n 1/k^2<2$ for all $n\ge1$, and watch the inductive step fail. Then prove the stronger statement $\sum_{k=1}^n 1/k^2\le 2-1/n$, and explain why the stronger statement was easier.
\end{exercise}

\begin{exercise}[$\star$]\label{ex:odds}
\emph{Nicomachus, proved.} Show that for $k\ge1$ the $k$ consecutive odd numbers $k^2-k+1$, $k^2-k+3$, \dots, $k^2+k-1$ add to $k^3$. Show that the blocks for $k=1,2,\dots,n$, placed end to end, are exactly the first $n(n+1)/2$ odd numbers. Deduce a third proof that $\sum_{k=1}^n k^3=\bigl(n(n+1)/2\bigr)^2$.
\end{exercise}

\begin{exercise}[$\star$]\label{ex:telprod}
For $n\ge2$, compute $\displaystyle\prod_{k=2}^n\left(1-\frac{1}{k^2}\right)$ and its limit as $n\to\infty$. Also compute $\displaystyle\sum_{k=1}^n\frac{k}{(k+1)!}$ and its limit.
\end{exercise}

\begin{exercise}[$\star\star$]\label{ex:pascalmod2}
\emph{Pascal's triangle modulo $2$.}
\begin{enumerate}[label=(\alph*)]
\item Show that $(1+x)^{2^m}\equiv1+x^{2^m}$ modulo $2$, in the sense that every coefficient of the difference is even. Use Exercise~\ref{ex:fermat}(a) and induction on $m$.
\item Deduce that row $2^m$ of Pascal's triangle is odd only at its two ends, and that row $2^m-1$ is entirely odd.
\end{enumerate}
\end{exercise}

\begin{exercise}[$\star\star$]\label{ex:fibsums}
With the Fibonacci numbers of Exercise~\ref{ex:fibgcd}, show that $\sum_{k=1}^n F_k=F_{n+2}-1$ and $\sum_{k=1}^n F_k^2=F_n F_{n+1}$. Draw the second identity: squares of sides $F_1,F_2,\dots,F_n$ tile an $F_n\times F_{n+1}$ rectangle.
\end{exercise}

\begin{remark}
The same recurrence has a closed form, used later when a numerical size is needed. Let $\varphi=(1+\sqrt5)/2$ and $\psi=(1-\sqrt5)/2$. Then
\[
F_n=\frac{\varphi^n-\psi^n}{\sqrt5}
\]
for every $n\ge0$. Both $\varphi$ and $\psi$ satisfy $t^2=t+1$, so $\varphi^n=\varphi^{n-1}+\varphi^{n-2}$ for $n\ge2$, and likewise for $\psi$. The sequence $G_n=(\varphi^n-\psi^n)/\sqrt5$ therefore satisfies the Fibonacci recurrence. Also $G_0=0$ and $G_1=(\varphi-\psi)/\sqrt5=1$. So $G_n=F_n$. Exercise~\ref{ex:fibmat} in Chapter~\ref{ch:linear} reads the same formula off a matrix.
\end{remark}

\begin{exercise}[$\star\star$]\label{ex:binomfrac}
Show that $\displaystyle\sum_{k=0}^n\frac{\binom nk}{k+1}=\frac{2^{n+1}-1}{n+1}$ and $\displaystyle\sum_{k=0}^n\frac{(-1)^k\binom nk}{k+1}=\frac{1}{n+1}$.
\end{exercise}

\begin{exercise}[$\star\star$]\label{ex:intpoly}
\emph{Integer-valued polynomials.} Let $P$ be a polynomial with complex coefficients such that $P(n)\in\Z$ for every $n\in\N$. Show that $P=\sum_k c_k\binom Xk$ with every $c_k\in\Z$. Deduce that $P(n)\in\Z$ for every $n\in\Z$, and that $d!\,P$ has integer coefficients if $\deg P=d$. Is $(X^3-X)/6$ integer-valued?
\end{exercise}

\begin{challenge}
\begin{enumerate}[label=(\arabic*)]
\item \emph{Faulhaber.} Compute $\sum_{k=1}^n k^5$ by Theorem~\ref{thm:newton}, and check that it is a polynomial in $N=n(n+1)/2$, namely $(4N^3-N^2)/3$. Faulhaber observed in $1631$ that every sum of odd powers $\sum_{k=1}^n k^{2m+1}$ is a polynomial in $N$. Why do the odd powers have this symmetry? Start from $S(n)=\sum_{k=1}^n k^{2m+1}$, extended as a polynomial, and from the relation $S(-1-n)=S(n)$.
\item \emph{Odd entries.} Count the odd entries in row $n$ of Pascal's triangle for $n\le16$, and guess the rule. It uses the binary expansion of $n$. Lucas' theorem, an exercise in the congruences chapter of Volume~2, proves the guess.
\item \emph{Squarefree central binomial coefficients.} For $n\ge1$, the numbers $\binom{2n}n$ begin $2,6,20,70,252,\dots$. Which of these are squarefree, meaning divisible by no square greater than $1$? Erd\H{o}s conjectured that $\binom{2n}n$ is never squarefree for $n>4$. Andrew Granville and Olivier Ramar\'e proved the conjecture in $1996$. Prove the easy part: if $n\ge1$ is not a power of $2$, then $4$ divides $\binom{2n}n$. Compare the number of factors $2$ in $(2n)!$ and in $(n!)^2$. The case $n=0$ is excluded: $\binom00=1$, and $0$ is not a power of $2$.
\end{enumerate}
Notes on the challenges are at the end of the chapter.
\end{challenge}

\begin{beauty}
Exercise~\ref{ex:fermat} is a first meeting of two worlds. A counting fact, the size of a binomial coefficient, proves an arithmetic fact, Fermat's congruence. Volume~2 carries that meeting forward, including a miniature form of the RSA encryption scheme. Exercise~\ref{ex:fibgcd} shows that the Fibonacci numbers form a divisibility sequence: a common divisor of indices divides the corresponding Fibonacci numbers in the same way. Lucas sequences and primality tests, later in Volume~2 and Volume~3, grow from that pattern.

The four proofs of $\sum k^3=(\sum k)^2$ are four ways of seeing one square: induction, a change of basis, a picture, and a double count. You will meet each of them again. The Catalan numbers of Exercise~\ref{ex:catalan}, in the next chapter, were studied by Euler in $1751$ while counting triangulations of a polygon. Volume~5 returns to them through generating functions.
\end{beauty}

The problems below connect this chapter with later tools. Problem~\ref{pb:millin}
uses the language of limits from Chapter~\ref{ch:sequences}.
Problem~\ref{pb:erdosheil} requires the polynomial method in
Problem~\ref{pb:nullstellensatz}, also in that chapter.
Problem~\ref{pb:bilinear} requires complex exponentials from
Chapter~\ref{ch:complex} and Cauchy--Schwarz from
Chapter~\ref{ch:inequalities}. You can read the questions now,
but those results are prerequisites for their given solutions.
They are not prerequisites for the counting chapter.

\begin{exercise}[$\star\star$]\label{pb:millin}
\begin{enumerate}[label=(\alph*)]
\item Show that $F_{2m}=F_m(F_{m-1}+F_{m+1})$ and $F_{2m-1}=F_m^2+F_{m-1}^2$ for $m\ge1$.
\item Show that $\sum_{n=0}^N 1/F_{2^n}=3-F_{2^N-1}/F_{2^N}$ for $N\ge1$. Deduce
\[
\sum_{n\ge0}\frac1{F_{2^n}}=1+1+\frac13+\frac1{21}+\frac1{987}+\cdots=\frac{7-\sqrt5}{2}.
\]
The language of this limit is in Chapter~\ref{ch:sequences}.
\end{enumerate}
\end{exercise}

\begin{exercise}[$\star\star\star$]\label{pb:erdosheil}
\emph{Sums of distinct elements.} For $A\subset\Z/p\Z$ let $A\dotplus A=\{a+a':a,a'\in A,\ a\ne a'\}$.
\begin{enumerate}[label=(\alph*)]
\item (Dias da Silva and Hamidoune, 1994; the argument below is due to Alon, Nathanson and Ruzsa, 1995.) Prove the Erd\H{o}s--Heilbronn conjecture: $|A\dotplus A|\ge\min(p,\,2|A|-3)$. Use Alon's combinatorial nullstellensatz, Problem~\ref{pb:nullstellensatz}(b) in Chapter~\ref{ch:sequences}.
\item Show that the bound is reached for every size $|A|\le(p+3)/2$.
\end{enumerate}
\end{exercise}

\begin{exercise}[$\star\star$]\label{pb:sauer}
\emph{The Sauer--Shelah lemma.} A family $\mathcal F$ of subsets of $[n]=\{1,\dots,n\}$ \emph{shatters} a set $T\subseteq[n]$ when $\{F\cap T:F\in\mathcal F\}$ is the set of all subsets of $T$. The \emph{VC-dimension} of $\mathcal F$ is the largest size of a shattered set.
\begin{enumerate}[label=(\alph*)]
\item Show that if $\mathcal F$ has VC-dimension at most $d$, then $|\mathcal F|\le\sum_{i=0}^d\binom ni$.
\item Show that the subsets of $n$ points on the real line cut out by intervals form a family of VC-dimension $2$ that attains the bound.
\end{enumerate}
\end{exercise}

\begin{exercise}[$\star\star\star$]\label{pb:bilinear}
\emph{Sums over boxes.} Let $p$ be an odd prime. Write $e(t)=\exp(2\pi i t)$. This problem uses two properties of that function: $e(t+1)=e(t)$, and $e(t)=1$ if and only if $t$ is an integer. The complex exponential is built in Chapter~\ref{ch:complex}. Parts (a) and (b) use Cauchy--Schwarz, Theorem~\ref{thm:cs} in Chapter~\ref{ch:inequalities}. Let $\mathbb F_p$ denote the residues modulo $p$.
\begin{enumerate}[label=(\alph*)]
\item Let $A,B\subseteq\mathbb F_p$ and let $r\not\equiv0$. Show that
\[
\Bigl|\sum_{a\in A}\sum_{b\in B}e\Bigl(\frac{rab}{p}\Bigr)\Bigr|\le\sqrt{p\,|A|\,|B|}.
\]
\item Let $U,V\subseteq\mathbb F_p^2$, let $x\not\equiv0$, and let $N(x)$ be the number of pairs $(u,v)\in U\times V$ with $u\cdot v=x$, the dot product $u_1v_1+u_2v_2$. Show that
\[
\Bigl|N(x)-\frac{|U|\,|V|}{p}\Bigr|\le\sqrt{p\,|U|\,|V|}.
\]
\item Let $A,B,C,D\subseteq\mathbb F_p$ with $|A|\,|B|\,|C|\,|D|>p^3$. Show that every $x\not\equiv0$ can be written $ab+cd$ with $a\in A$, $b\in B$, $c\in C$ and $d\in D$. In particular, if $|A|>p^{3/4}$, every nonzero residue has the form $ab+cd$ with $a,b,c,d\in A$.
\item Let $A=B=\{1,\dots,m\}$ with $m\le\sqrt p/3$. Show that the sum in (a), with $r=1$, has real part at least $m^2/2$. Taking $m=\lfloor\sqrt p/3\rfloor$, conclude that (a) cannot be improved by more than a constant factor.
\end{enumerate}
\end{exercise}

\section{Hints}

\begin{hint}{ex:cubes}
In the inductive step, factor $(n+1)^2$ out of $\bigl(n(n+1)/2\bigr)^2+(n+1)^3$ before you expand.
\end{hint}
\begin{hint}{ex:tel3}
Compute $\dfrac1{k(k+1)}-\dfrac1{(k+1)(k+2)}$, and compare it with $\dfrac1{k(k+1)(k+2)}$.
\end{hint}
\begin{hint}{ex:kbinom}
The term $k\binom nk$ counts committees of size $k$ with a chosen chair. Try choosing the chair first. For the algebraic proof, cancel factorials before you sum.
\end{hint}
\begin{hint}{ex:hockey}
Induct on $n$, using Pascal's rule. Or sort the $(r+1)$-element subsets of $\{1,\dots,n+1\}$ by their largest element.
\end{hint}
\begin{hint}{ex:squares}
Write $\binom nk^2=\binom nk\binom n{n-k}$ and use Vandermonde. For the inequality, expand $\sum_{0\le i<j\le n}(b_i-b_j)^2$ and use that the sum is nonnegative.
\end{hint}
\begin{hint}{ex:fermat}
(a) From formula (b), $k!\,(p-k)!\binom pk=p!$. The prime $p$ divides the right side and divides no factor of $k!\,(p-k)!$. (b) Induct on $a$, and expand $(a+1)^p$.
\end{hint}
\begin{hint}{ex:fibgcd}
(a) Fix $m$ and induct on $n$. You need two base cases, because the recurrence looks back two steps. (b) A common divisor of $F_n$ and $F_{n+1}$ divides $F_{n-1}$. (c) Part (a) gives $\gcd(F_{m+n},F_n)=\gcd(F_m,F_n)$. That is Euclid's algorithm, applied to the indices.
\end{hint}
\begin{hint}{ex:cubesgeo}
The region $L_k$ is the row segment $k\cdot1,\dots,k\cdot k$ together with the column segment $1\cdot k,\dots,(k-1)\cdot k$. See the right half of Figure~\ref{fig:pascal}.
\end{hint}
\begin{hint}{ex:harmonic}
Let $2^s$ be the largest power of $2$ with $2^s\le n$. In the sum $H_n$, only the term $1/2^s$ has $2^s$ in the denominator.
\end{hint}
\begin{hint}{ex:sumsq}
From the weak bound $S_n<2$ you only get $S_{n+1}<2+1/(n+1)^2$, which does not stay below $2$. For the stronger bound, check that $1/n-1/(n+1)^2\ge1/(n+1)$.
\end{hint}
\begin{hint}{ex:odds}
The $k$ numbers form an arithmetic progression of odd numbers. Their average is $k^2$. For the second claim, check that block $k+1$ begins at the next odd number after block $k$.
\end{hint}
\begin{hint}{ex:telprod}
Factor $1-1/k^2=(k-1)(k+1)/k^2$. That is two telescoping products. For the sum, write $k/(k+1)!=1/k!-1/(k+1)!$.
\end{hint}
\begin{hint}{ex:pascalmod2}
(a) $(1+x)^{2^{m+1}}=\bigl((1+x)^{2^m}\bigr)^2$. If $A-B$ has even coefficients, so does $A^2-B^2$. (b) For the row $2^m-1$, Pascal's rule says that $\binom{2^m}{k}=\binom{2^m-1}{k-1}+\binom{2^m-1}{k}$ is even for $0<k<2^m$. Walk along the row from $k=0$.
\end{hint}
\begin{hint}{ex:fibsums}
Both sums telescope: $F_k=F_{k+2}-F_{k+1}$, and $F_k^2=F_k(F_{k+1}-F_{k-1})=F_k F_{k+1}-F_{k-1}F_k$.
\end{hint}
\begin{hint}{ex:binomfrac}
Check from formula (b) that $\binom nk/(k+1)=\binom{n+1}{k+1}/(n+1)$. The sums $\sum_k\binom{n+1}{k+1}$ and $\sum_k(-1)^k\binom{n+1}{k+1}$ are then pieces of the binomial theorem.
\end{hint}
\begin{hint}{ex:intpoly}
Theorem~\ref{thm:newton} gives $c_k=(\Delta^k P)(0)$, an integer combination of $P(0),\dots,P(k)$. For a negative integer, $\binom{-n}{k}=(-1)^k\binom{n+k-1}{k}$.
\end{hint}

\begin{hint}{pb:millin}
Use the addition formula from Exercise~\ref{ex:fibgcd}. For the fraction in (b), you also need Cassini's identity $F_{m-1}F_{m+1}-F_m^2=(-1)^m$. It has a short induction from the Fibonacci recurrence; the solution writes it out. The matrix proof is Exercise~\ref{ex:fibmat}.
\end{hint}

\begin{hint}{pb:erdosheil}
Start from the Cauchy--Davenport polynomial of Problem~\ref{pb:cauchydav}, and multiply by $x-y$ to remove the diagonal $a=a'$. The coefficient you need is a difference of two binomial coefficients, hence a Catalan number. Apply Problem~\ref{pb:nullstellensatz}.
\end{hint}

\begin{hint}{pb:sauer}
(a) Induct on $n$. Compare $\mathcal F$ with the family of sets $F\setminus\{n\}$, and with the sets $F$ for which both $F$ and $F\cup\{n\}$ lie in $\mathcal F$.
\end{hint}

\begin{hint}{pb:bilinear}
(a) Apply Cauchy--Schwarz (Theorem~\ref{thm:cs}) in the variable $a$, then let $a$ run through all of $\mathbb F_p$. The complete exponential sum is a geometric series. (b) Write $N(x)$ by summing the geometric character over the residue $u\cdot v-x$. Isolate the term $r=0$. (c) Take $U=A\times C$ and $V=B\times D$. (d) Estimate the cosines.
\end{hint}

\section{Solutions}

\begin{solution}{ex:cubes}
\idea The first proof in Section~\ref{sec:fourproofs} adds one cube to the inductive square. Factor $(n+1)^2$ out before you expand, so the next square appears.
\writeup Write that proof. A common slip is to multiply out both squares first. The factor $(n+1)^2$ is already in the hypothesis, and the new cube is $(n+1)^2\cdot(n+1)$, so pull that factor out and the algebra is one line. Check $n=2$: $1+8=9$ and $(2\cdot3/2)^2=9$.
\end{solution}

\begin{solution}{ex:tel3}
\idea Proposition~\ref{prop:telescope} asks for a difference. The model $1/(k(k+1))=1/k-1/(k+1)$ suggests trying one level higher.
\writeup Compute
\[
\frac1{k(k+1)}-\frac1{(k+1)(k+2)}=\frac{(k+2)-k}{k(k+1)(k+2)}=\frac{2}{k(k+1)(k+2)}.
\]
So $1/(k(k+1)(k+2))$ is half of that difference. Summing from $k=1$ to $n$ cancels every intermediate term:
\[
S_n=\frac12\left(\frac1{1\cdot2}-\frac1{(n+1)(n+2)}\right)=\frac14-\frac1{2(n+1)(n+2)}.
\]
As $n\to\infty$, the second term tends to $0$, so $S_n\to1/4$. The finite formula already shows the error: it is smaller than $1/(2n^2)$.
\end{solution}

\begin{solution}{ex:kbinom}
\idea The factor $k$ marks one element of a $k$-element set. Call that element the chair of a committee. Choosing the chair first leaves a free choice for every other person.
\writeup Let $n\ge1$. Count pairs (committee, chair) among $n$ people. Choosing the committee size $k$, then the committee, then the chair, gives $\sum_{k=0}^n k\binom nk$. Choosing the chair first gives $n$ choices, and then any subset of the remaining $n-1$ people may join, in $2^{n-1}$ ways. The two counts are equal, so the sum is $n\,2^{n-1}$.

The same identity is algebraic. For $1\le k\le n$,
\[
k\binom nk=k\cdot\frac{n!}{k!\,(n-k)!}=n\cdot\frac{(n-1)!}{(k-1)!\,(n-k)!}=n\binom{n-1}{k-1}.
\]
Sum over $k$. The binomial theorem gives $\sum_{k\ge1}\binom{n-1}{k-1}=2^{n-1}$. For $n=0$ the original sum is $0$, because the only subset is empty and has no chair.

Once differentiation is available, differentiate $(1+x)^n=\sum_k\binom nk x^k$, multiply by $x$, and set $x=1$. For $n\ge1$ you recover the same identity.
\end{solution}

\begin{solution}{ex:hockey}
\idea When $n$ grows by one, the left side gains one term. Pascal's rule says the right side gains that same term. For a counting proof, sort subsets by their largest element.
\writeup Induct on $n$, for a fixed $r$, starting at $n=r$. Both sides equal $1$ when $n=r$, because the sum is the single term $\binom rr=1$ and $\binom{r+1}{r+1}=1$. If the identity holds up through $n$, Pascal's rule gives
\[
\binom{n+2}{r+1}=\binom{n+1}{r+1}+\binom{n+1}{r}=\sum_{k=r}^n\binom kr+\binom{n+1}{r}=\sum_{k=r}^{n+1}\binom kr.
\]
That is the step from $n$ to $n+1$.

Counting proof: an $(r+1)$-element subset of $\{1,\dots,n+1\}$ has a largest element $k+1$, with $r\le k\le n$. The rest is an $r$-element subset of $\{1,\dots,k\}$. Each such subset arises once.

For the cubes, follow the second proof in Section~\ref{sec:fourproofs}. The slip is to stop the hockey stick at $\binom n{r+1}$. Summing $\binom kr$ from $k=r$ through $n$ lands on $\binom{n+1}{r+1}$, one row past the last term. For $n=3$ that form is $\binom42+6\binom43+6\binom44=36$.
\end{solution}

\begin{solution}{ex:squares}
\idea Vandermonde needs two lower indices that add to a constant. The symmetry $\binom nk=\binom n{n-k}$ supplies them. For the bound, the only information we use is that the row sums to $2^n$. Squared gaps between the entries measure how far they are from being equal.
\writeup Proposition~\ref{prop:vandermonde} with $m=n$ and $r=n$ gives
\[
\binom{2n}{n}=\sum_{k=0}^n\binom nk\binom n{n-k}=\sum_{k=0}^n\binom nk^2.
\]
For the inequality, take any real numbers $b_0,\dots,b_n$. Expand
\begin{align*}
\sum_{0\le i<j\le n}(b_i-b_j)^2
&=\sum_{0\le i<j\le n}(b_i^2+b_j^2-2b_ib_j).
\end{align*}
There are $n+1$ indices, so each fixed index meets $n$ others. Each $b_t^2$ therefore appears $n$ times in the sum of squares, and
\[
\sum_{0\le i<j\le n}(b_i^2+b_j^2)=n\sum_{t=0}^n b_t^2.
\]
Also $\bigl(\sum_t b_t\bigr)^2=\sum_t b_t^2+2\sum_{i<j}b_ib_j$. Combine these:
\[
(n+1)\sum_{t=0}^n b_t^2-\Bigl(\sum_{t=0}^n b_t\Bigr)^2=\sum_{0\le i<j\le n}(b_i-b_j)^2\ge0.
\]
Apply this to $b_k=\binom nk$. The binomial theorem gives $\sum b_k=2^n$, so $4^n\le(n+1)\binom{2n}{n}$. For $n=0$ both sides are $1$. Chapter~\ref{ch:inequalities} meets the same comparison again as Cauchy--Schwarz, Theorem~\ref{thm:cs}.
\end{solution}

\begin{solution}{ex:fermat}
\idea In $\binom pk=p!/(k!(p-k)!)$, the prime $p$ stands in the numerator, and every factor in the denominator is smaller than $p$. Euclid's lemma is what turns ``the $p$ cannot cancel'' into a proof. Part (b) is then an induction, because $(a+1)^p$ and $a^p+1$ agree modulo $p$.
\writeup (a) By Proposition~\ref{prop:binom}(b), $k!\,(p-k)!\binom pk=p!$. The right side is divisible by $p$. For $0<k<p$, the integer $k!\,(p-k)!$ is a product of integers from $1$ to $p-1$. None of those factors is divisible by $p$. Euclid's lemma, proved earlier in this chapter, says that a prime dividing a product divides one of the factors. The contrapositive says that $p$ does not divide $k!\,(p-k)!$. That factorial product is therefore coprime to $p$. Coprime cancellation, from the same proposition, yields $p\mid\binom pk$.

(b) Induct on $a$. For $a=0$ both sides are $0$. If $a^p\equiv a\pmod p$, expand $(a+1)^p$ by the binomial theorem. Part (a) says that every coefficient $\binom pk$ with $0<k<p$ is divisible by $p$, so
\[
(a+1)^p=\sum_{k=0}^p\binom pk a^k\equiv a^p+1\equiv a+1\pmod p.
\]
\end{solution}

\begin{solution}{ex:fibgcd}
\idea The recurrence looks two steps back, so an induction on $n$ needs two base cases and, in the step, two earlier cases. A common divisor of two consecutive Fibonacci numbers divides their difference, which is the previous one. The gcd of the values then follows Euclid's algorithm on the indices.
\writeup (a) Fix $m\ge1$, and induct on $n$. For $n=0$, the claim is $F_m=F_m F_1+F_{m-1}F_0$, and $F_1=1$, $F_0=0$. For $n=1$, the claim is $F_{m+1}=F_m F_2+F_{m-1}F_1=F_m+F_{m-1}$, which is the recurrence. If the claim holds for $n$ and for $n+1$, add those two equations. The recurrence on each column produces the claim for $n+2$.

(b) For $n=0$, $\gcd(F_0,F_1)=\gcd(0,1)=1$. For $n\ge1$, a common divisor of $F_n$ and $F_{n+1}$ divides $F_{n+1}-F_n=F_{n-1}$. Repeating, it divides $F_1=1$.

(c) If $n=0$ and $m>0$, then $\gcd(F_m,F_0)=\gcd(F_m,0)=F_m=F_{\gcd(m,0)}$. If $m=n>0$, then $\gcd(F_m,F_m)=F_m=F_{\gcd(m,m)}$. The identity is symmetric in $m$ and $n$, so we may assume $m>n\ge1$. Part (a), with the indices swapped in the sense $F_m=F_{(m-n)+n}$, gives
\[
F_m=F_n F_{m-n+1}+F_{n-1}F_{m-n}.
\]
Hence $\gcd(F_m,F_n)=\gcd(F_{n-1}F_{m-n},F_n)$. By (b), $\gcd(F_{n-1},F_n)=1$. Any common divisor of $F_{n-1}F_{m-n}$ and $F_n$ is therefore coprime to $F_{n-1}$, so Euclid's lemma puts it as a divisor of $F_{m-n}$. Conversely, every common divisor of $F_{m-n}$ and $F_n$ divides the product. So
\[
\gcd(F_m,F_n)=\gcd(F_{m-n},F_n).
\]
The indices obey the same subtraction as the ordinary gcd. Repeating reaches $\gcd(F_d,F_0)=F_d$ with $d=\gcd(m,n)>0$.
\end{solution}

\begin{solution}{ex:cubesgeo}
\idea The fourth proof in Section~\ref{sec:fourproofs} cuts the table into the regions $L_k$. Each region is two arithmetic progressions with factor $k$, and the table total is already $\bigl(\sum i\bigr)^2$.
\writeup Write that proof. The usual mistake is to list the corner $(k,k)$ in the row and again in the column. That counts $k^2$ twice, and the region sums to $k^3+k^2$. Keep the corner in the row only. For $n=1$ the table is the single cell $1$, and both sides equal $1$.
\end{solution}

\begin{solution}{ex:harmonic}
\idea To show that a rational number is not an integer, keep one prime in the denominator after clearing the others. Here the prime is $2$, and the special term is $1/2^s$.
\writeup Let $n\ge2$, and let $2^s$ be the largest power of $2$ with $2^s\le n$. Then $s\ge1$. Among $1,\dots,n$, the only multiple of $2^s$ is $2^s$ itself: the next multiple is $2^{s+1}>n$. Let $L=\operatorname{lcm}(1,\dots,n)$. The exponent of $2$ in $L$ is exactly $s$.

Consider $L H_n=\sum_{k=1}^n L/k$. For $k\ne2^s$, the power of $2$ dividing $k$ is at most $2^{s-1}$, so $L/k$ is even. The remaining term $L/2^s$ is odd. A sum of one odd integer and some even integers is odd. If $H_n$ were an integer, then $L H_n$ would be even, because $L$ is even. That is a contradiction. So $H_n$ is not an integer.

For a check, $H_2=3/2$ and $H_3=11/6$, neither an integer. The argument does not say how large the denominator is, only that a factor $2$ survives.
\end{solution}

\begin{solution}{ex:sumsq}
\idea The bound $S_n<2$ leaves no room. Adding $1/(n+1)^2$ may step over $2$, and the inductive hypothesis cannot forbid it. A bound $S_n\le2-1/n$ leaves a gap of size $1/n$, and that gap is larger than the next term.
\writeup Let $S_n=\sum_{k=1}^n 1/k^2$, and let $P(n)$ be the statement $S_n\le2-1/n$. For $n=1$, $S_1=1$ and $2-1/1=1$. Suppose $P(n)$ holds. Then
\[
S_{n+1}\le2-\frac1n+\frac1{(n+1)^2}.
\]
We need this to be at most $2-1/(n+1)$. That rearrangement is
\[
\frac1n-\frac1{(n+1)^2}\ge\frac1{n+1},
\]
or equivalently $1/n-1/(n+1)\ge1/(n+1)^2$. The left side equals $1/(n(n+1))$. The denominator $n(n+1)$ is smaller than $(n+1)^2$, because $n<n+1$. A smaller positive denominator gives a larger fraction, so $1/(n(n+1))>1/(n+1)^2$. The gap is large enough, and $P(n+1)$ holds. In particular $S_n\le2-1/n<2$.

The weaker statement was harder because its hypothesis, $S_n<2$, does not know how much room remains under $2$. The stronger statement is a stronger hypothesis, and the step spends that extra information. Chapter~\ref{ch:sequences} will show that the infinite sum equals $\pi^2/6$, about $1.645$, which sits strictly under $2$.
\end{solution}

\begin{solution}{ex:odds}
\idea The third proof in Section~\ref{sec:fourproofs} groups the odd numbers into blocks of length $k$ that add to $k^3$.
\writeup Write that proof. The slip is a gap or an overlap between blocks. Block $k$ ends at $k^2+k-1$, and block $k+1$ begins at $k^2+k+1$, the next odd number. For $k=3$ the block is $7+9+11=27$.
\end{solution}

\begin{solution}{ex:telprod}
\idea Split $1-1/k^2$ into two fractions, each of which cancels in its own product. Split the numerator $k$ so that the summand is a difference of reciprocal factorials.
\writeup For $n\ge2$,
\[
1-\frac1{k^2}=\frac{(k-1)(k+1)}{k^2}=\frac{k-1}{k}\cdot\frac{k+1}{k}.
\]
The product is
\[
\prod_{k=2}^n\frac{k-1}{k}\cdot\prod_{k=2}^n\frac{k+1}{k}=\frac1n\cdot\frac{n+1}{2}=\frac{n+1}{2n}.
\]
As $n\to\infty$, this tends to $1/2$.

For the sum, $k=(k+1)-1$, so
\[
\frac{k}{(k+1)!}=\frac{k+1}{(k+1)!}-\frac1{(k+1)!}=\frac1{k!}-\frac1{(k+1)!}.
\]
Summing from $k=1$ to $n$ leaves $1/1!-1/(n+1)!=1-1/(n+1)!$, which tends to $1$.
\end{solution}

\begin{solution}{ex:pascalmod2}
\idea Modulo $2$, the freshman's dream $(a+b)^2\equiv a^2+b^2$ holds, because the cross term $2ab$ is even. That is Exercise~\ref{ex:fermat}(a) at $p=2$. Squaring repeatedly produces the powers $2^m$.
\writeup (a) For $m=0$, $(1+x)^1-(1+x)=0$, so there is nothing to prove. Suppose $(1+x)^{2^m}=1+x^{2^m}+2R(x)$ for some polynomial $R$ with integer coefficients. Square both sides:
\begin{align*}
(1+x)^{2^{m+1}}
&=(1+x^{2^m})^2+4(1+x^{2^m})R(x)+4R(x)^2\\
&=1+2x^{2^m}+x^{2^{m+1}}+4(\cdots).
\end{align*}
Every term except $1$ and $x^{2^{m+1}}$ has an even coefficient.

(b) The coefficients of $(1+x)^{2^m}$ are row $2^m$ of Pascal's triangle. By (a), $\binom{2^m}{k}$ is even for $0<k<2^m$, and the two ends $\binom{2^m}{0}$ and $\binom{2^m}{2^m}$ equal $1$. For the previous row, Pascal's rule says
\[
\binom{2^m}{k}=\binom{2^m-1}{k-1}+\binom{2^m-1}{k}.
\]
For $0<k<2^m$ the left side is even, so the two entries of row $2^m-1$ have the same parity. Starting from $\binom{2^m-1}{0}=1$, which is odd, every entry of that row is odd.
\end{solution}

\begin{solution}{ex:fibsums}
\idea The recurrence says $F_k=F_{k+2}-F_{k+1}$, a difference. For the squares, rewrite one factor $F_k$ as $F_{k+1}-F_{k-1}$. The product becomes a difference of consecutive products $F_k F_{k+1}$.
\writeup For $k\ge1$, $F_{k+2}-F_{k+1}=F_k$. Sum from $k=1$ to $n$:
\[
\sum_{k=1}^n F_k=\sum_{k=1}^n(F_{k+2}-F_{k+1})=F_{n+2}-F_2=F_{n+2}-1,
\]
because $F_2=1$. For the squares, and for $k\ge1$,
\[
F_k^2=F_k(F_{k+1}-F_{k-1})=F_k F_{k+1}-F_{k-1}F_k.
\]
The $k=1$ term uses $F_0=0$. Summing leaves $F_n F_{n+1}-F_0 F_1=F_n F_{n+1}$.

Figure~\ref{fig:fibrect} draws the case $n=5$. The squares of sides $F_1,\dots,F_5$ fill an $F_5\times F_6$ rectangle, so the area $F_5 F_6$ equals the sum of the five square areas. The same placement, stopped at $n$, fills an $F_n\times F_{n+1}$ rectangle.
\end{solution}

\begin{figure}[ht]
\centering
\fig{ch01-fibrect}
\caption{Squares of sides $F_1=1$, $F_2=1$, $F_3=2$, $F_4=3$ and $F_5=5$ tile a rectangle of size $5\times8$, that is $F_5\times F_6$. The area says $\sum_{k=1}^5 F_k^2=F_5 F_6$.}
\label{fig:fibrect}
\end{figure}

\begin{solution}{ex:binomfrac}
\idea A factor $k+1$ in the denominator, next to $\binom nk$, can be absorbed into a binomial coefficient in the next row. Formula (b) checks the absorption. The sums then become partial binomial expansions.
\writeup For $0\le k\le n$,
\[
\frac{\binom nk}{k+1}=\frac{n!}{(k+1)!\,(n-k)!}=\frac{1}{n+1}\cdot\frac{(n+1)!}{(k+1)!\,(n-k)!}=\frac{1}{n+1}\binom{n+1}{k+1}.
\]
Sum over $k$. The change of index $j=k+1$ runs from $1$ to $n+1$, so
\[
\sum_{k=0}^n\frac{\binom nk}{k+1}=\frac1{n+1}\sum_{j=1}^{n+1}\binom{n+1}{j}=\frac{2^{n+1}-1}{n+1},
\]
because the full sum of row $n+1$ is $2^{n+1}$ and the missing $j=0$ term is $1$. With the signs,
\begin{align*}
\sum_{k=0}^n\frac{(-1)^k\binom nk}{k+1}
&=\frac1{n+1}\sum_{j=1}^{n+1}(-1)^{j-1}\binom{n+1}{j}.
\end{align*}
The binomial theorem gives $\sum_{j=0}^{n+1}(-1)^j\binom{n+1}{j}=0$ when $n+1\ge1$. The $j=0$ term is $1$, so $\sum_{j\ge1}(-1)^j\binom{n+1}{j}=-1$. Multiplying by $-1$ produces $\sum_{j\ge1}(-1)^{j-1}\binom{n+1}{j}=1$. The signed sum is therefore $1/(n+1)$.
\end{solution}

\begin{solution}{ex:intpoly}
\idea Newton's formula writes the coefficients $c_k$ as $(\Delta^k P)(0)$. If the values $P(0),P(1),\dots$ are integers, then so are their differences. Once the $c_k$ are integers, the binomial polynomials take integer values at every integer, positive or negative.
\writeup If $P=0$, take every $c_k=0$. Otherwise let $d=\deg P$. The basis argument in Theorem~\ref{thm:newton} uses only degrees and leading coefficients, so it applies to complex coefficients exactly as stated in the chapter. Write $P=\sum_{k\le d}c_k\binom Xk$. Then $c_k=(\Delta^k P)(0)$. Each difference of integer values is an integer, and $P(0),\dots,P(k)$ are integers, so $c_k\in\Z$. Explicitly, by a separate induction on $k$,
\[
(\Delta^k P)(0)=\sum_{j=0}^k(-1)^{k-j}\binom kj P(j).
\]

For $n\in\Z$ and $k\ge0$, the number $\binom nk=n(n-1)\cdots(n-k+1)/k!$ is an integer. For $n\ge0$ it counts subsets, or it is zero when $n<k$. For $n=-m$ with $m>0$,
\[
\binom{-m}{k}=(-1)^k\frac{m(m+1)\cdots(m+k-1)}{k!}=(-1)^k\binom{m+k-1}{k},
\]
which is an integer. So $P(n)\in\Z$ for every integer $n$.

Finally, for $k\le d$ the polynomial $d!\binom Xk$ has integer coefficients, because $k!$ divides $d!$ and the numerator $X(X-1)\cdots(X-k+1)$ is a monic polynomial with integer coefficients. Hence $d!\,P$ lies in $\Z[X]$.

The polynomial $(X^3-X)/6$ equals $\binom{X+1}{3}$. Indeed
\[
\binom{X+1}{3}=\frac{(X+1)X(X-1)}{6}=\frac{X^3-X}{6}.
\]
It is therefore integer-valued. Among any three consecutive integers, one is divisible by $2$ and one is divisible by $3$, so the product $(X-1)X(X+1)$ is divisible by $6$. The difference table at $0,1,2,3$ begins $0,0,1,4$, then $0,1,3$, then $1,2$, then $1$, which is the same conclusion from Theorem~\ref{thm:newton}: $(X^3-X)/6=\binom X3+\binom X2$.
\end{solution}

\subsection*{Notes on the challenges}

(1) The difference table of $0^5,1^5,2^5,3^5,4^5,5^5$, that is $0,1,32,243,1024,3125$, begins with the first entries $0,1,30,150,240,120$. Theorem~\ref{thm:newton} therefore gives
\[
k^5=\binom k1+30\binom k2+150\binom k3+240\binom k4+120\binom k5.
\]
Summing by Corollary~\ref{cor:sumpoly}, and shifting the upper limit from $n-1$ to $n$, gives
\[
\sum_{k=1}^n k^5=\binom{n+1}{2}+30\binom{n+1}{3}+150\binom{n+1}{4}+240\binom{n+1}{5}+120\binom{n+1}{6}.
\]
Expanding and setting $N=n(n+1)/2$ produces
\[
\sum_{k=1}^n k^5=\frac{n^2(n+1)^2(2n^2+2n-1)}{12}=\frac{4N^3-N^2}{3}.
\]
Both sides agree at $n=1$ and at $n=2$, which is a useful check of the arithmetic: the sum is $1$, then $33$.

Here is the symmetry for odd powers. By the corollary, $S(n)=\sum_{k=1}^n k^{2m+1}$ agrees on the nonnegative integers with a polynomial, still called $S$, satisfying $S(0)=0$ and $S(n+1)-S(n)=(n+1)^{2m+1}$ as an identity of polynomials. Set $n=-1$ in the difference: $S(0)-S(-1)=0$, so $S(-1)=0$.

Let $T(n)=S(-n-1)$. Then
\begin{align*}
T(n+1)-T(n)
&=S(-n-2)-S(-n-1)\\
&=-(S(-n-1)-S(-n-2)).
\end{align*}
The forward difference of $S$ at $-n-2$ is $(-n-1)^{2m+1}$. The exponent $2m+1$ is odd, so
\[
-(-n-1)^{2m+1}=(n+1)^{2m+1}.
\]
Thus $T$ and $S$ have the same forward difference. Their difference is constant. At $0$, that constant is $T(0)-S(0)=S(-1)=0$. So $S(-n-1)=S(n)$.

In the coordinate $z=n+1/2$, the identity $S(-n-1)=S(n)$ says that $S$ is an even function of $z$. An even polynomial in $z$ is a polynomial in $z^2$. And
\[
z^2=\Bigl(n+\frac12\Bigr)^2=n(n+1)+\frac14=2N+\frac14,
\]
so a polynomial in $z^2$ is a polynomial in $N$. This uses that the exponent is odd. For an even power the sign would not flip, and the same conclusion fails.

Faulhaber's published formulas do not include their derivation. Knuth's article, listed in the sources, examines what the $1631$ book contains and which methods were available.

(2) The number of odd entries in row $n$ is $2^{s(n)}$, where $s(n)$ is the number of ones in the binary expansion of $n$. This count is due to Glaisher ($1899$). Row $7=111_2$ has $8$ odd entries. Row $8=1000_2$ has $2$. Lucas' theorem, an exercise in the congruences chapter of Volume~2, says that $\binom nk$ is odd exactly when every binary digit of $k$ is less than or equal to the corresponding digit of $n$. The count of such $k$ is the product, over the binary digits of $n$, of one more than that digit, which is $2^{s(n)}$.

(3) For $n\ge1$, $\binom{2n}{n}=2\binom{2n-1}{n-1}$ is even. Legendre's formula for the exponent of a prime in $n!$ gives the exact power of $2$: it equals $s(n)$, the number of ones in the binary expansion of $n$. Volume~2 derives that exponent from Legendre's formula, and Volume~5 records Kummer's theorem, which counts the same exponent by the carries when $n$ is added to $n$ in base $2$. Thus for $n\ge1$ one has $4\mid\binom{2n}{n}$ unless $s(n)=1$, that is, unless $n$ is a power of $2$. For $n=2^m$ an odd square has to be found in the factorisation, and that is the hard part of the theorem of Granville and Ramar\'e. The squarefree values for $n\ge1$ are exactly $\binom21=2$, $\binom42=6$ and $\binom84=70$.

\begin{solution}{pb:millin}
\idea The denominators $F_{2^n}$ double their indices at every step, which suggests a difference $F_{m-1}/F_m-F_{2m-1}/F_{2m}$. The first partial sums $2$, $7/3$, $50/21$ sit just under $3$, and that suggests the closed form in (b).
\writeup (a) Exercise~\ref{ex:fibgcd}(a) with $n=m$ gives $F_{2m}=F_m F_{m+1}+F_{m-1}F_m=F_m(F_{m-1}+F_{m+1})$. With $n=m-1$, and with $m\ge1$, it gives $F_{2m-1}=F_m F_m+F_{m-1}F_{m-1}$.

(b) First we record Cassini's identity $F_{m-1}F_{m+1}-F_m^2=(-1)^m$ for $m\ge1$, by induction on $m$. For $m=1$, $F_0 F_2-F_1^2=-1$. If the identity holds for $m$, then
\begin{align*}
F_m F_{m+2}-F_{m+1}^2
&=F_m(F_{m+1}+F_m)-F_{m+1}(F_m+F_{m-1})\\
&=F_m^2-F_{m-1}F_{m+1}\\
&=-(-1)^m=(-1)^{m+1}.
\end{align*}
The matrix proof of the same identity is Exercise~\ref{ex:fibmat}.

Now let $m\ge2$ be even, so $(-1)^m=1$. Then
\begin{align*}
\frac{F_{m-1}}{F_m}-\frac{F_{2m-1}}{F_{2m}}
&=\frac{F_{m-1}F_m(F_{m-1}+F_{m+1})-F_m(F_m^2+F_{m-1}^2)}{F_m F_{2m}}\\
&=\frac{F_{m-1}F_{m+1}-F_m^2}{F_{2m}}\\
&=\frac{1}{F_{2m}}.
\end{align*}
Set $m=2^N$ with $N\ge1$. The identity says
\[
\Bigl(3-\frac{F_{2^{N+1}-1}}{F_{2^{N+1}}}\Bigr)-\Bigl(3-\frac{F_{2^N-1}}{F_{2^N}}\Bigr)=\frac1{F_{2^{N+1}}}.
\]
For $N=1$, $1/F_1+1/F_2=1+1=2$, and $3-F_1/F_2=3-1=2$. Adding the differences from $N=1$ up through $N-1$ gives the finite sum in the statement.

For the infinite sum, use the closed form in the remark after Exercise~\ref{ex:fibsums}. Since $|\psi|<\varphi$, the powers $(\psi/\varphi)^n$ tend to $0$. Therefore
\[
\frac{F_{n-1}}{F_n}=\frac1\varphi\cdot\frac{1-(\psi/\varphi)^{n-1}}{1-(\psi/\varphi)^n}\to\frac1\varphi=\frac{\sqrt5-1}{2}
\]
as $n\to\infty$. The language of this limit is in Chapter~\ref{ch:sequences}. Hence the partial sums tend to $3-(\sqrt5-1)/2=(7-\sqrt5)/2$.

\tools Exercise~\ref{ex:fibgcd}; Cassini, proved above and again in Exercise~\ref{ex:fibmat}.
\leads This is Millin's series ($1974$). The sum $\sum_{n\ge1}1/F_n$ has no known closed form. Andr\'e-Jeannin proved in $1989$ that it is irrational.
\end{solution}

\begin{solution}{pb:erdosheil}
\idea Multiply the Cauchy--Davenport polynomial by $x-y$. That factor vanishes on the diagonal $a=a'$, which is exactly what ``distinct summands'' removes. The coefficient that the nullstellensatz needs becomes a Catalan number.
\writeup (a) Let $k=|A|$. For $k\le2$ the lower bound $\min(p,2k-3)$ is at most $1$. A set of size $0$ or $1$ has no sum of two distinct elements, and the bound is negative, so the inequality holds. A set of size $2$ has exactly one such sum, and the bound equals $1$. The rest of the argument takes $k\ge3$. If $2k-3\ge p$, then $p$ is odd, because $p=2$ forces $k\le2$. The set $A$ then contains a subset $A'$ with $|A'|=(p+3)/2$. It is enough to prove $|A'\dotplus A'|\ge p$, so we may assume $2k-3\le p$, and therefore $2k-4<p$.

Suppose, for a contradiction, that $|A\dotplus A|\le2k-4$. Choose a set $C$ with $A\dotplus A\subseteq C$ and $|C|=2k-4$. The polynomial
\[
f(x,y)=(x-y)\prod_{c\in C}(x+y-c)
\]
vanishes on $A\times A$: on the diagonal because of the factor $x-y$, and off the diagonal because $x+y\in C$. Its total degree is $2k-3=(k-1)+(k-2)$. The coefficient of $x^{k-1}y^{k-2}$ in $(x-y)(x+y)^{2k-4}$ comes from two contributions. The factor $x$ times the term $\binom{2k-4}{k-2}x^{k-2}y^{k-2}$ contributes $\binom{2k-4}{k-2}$. The factor $-y$ times the term $\binom{2k-4}{k-1}x^{k-1}y^{k-3}$ contributes $-\binom{2k-4}{k-1}$. The ratio of those binomial coefficients is $(k-2)/(k-1)$, so the coefficient equals
\[
\binom{2k-4}{k-2}\Bigl(1-\frac{k-2}{k-1}\Bigr)=\frac1{k-1}\binom{2k-4}{k-2},
\]
the Catalan number $C_{k-2}$. Since $2k-4<p$, the prime $p$ does not divide $\binom{2k-4}{k-2}$, nor the denominator $k-1<p$. Problem~\ref{pb:nullstellensatz}(b), in Chapter~\ref{ch:sequences}, with $S_1=S_2=A$, $t_1=k-1$ and $t_2=k-2$, supplies a point of $A\times A$ where $f\ne0$. That contradicts the vanishing. So $|A\dotplus A|\ge\min(p,2|A|-3)$.

(b) For $A=\{0,1,\dots,k-1\}$, the sums of two distinct elements are $1,2,\dots,2k-3$. These are $2k-3$ distinct residues when $2k-3\le p$, that is when $k\le(p+3)/2$.

\tools Problems~\ref{pb:nullstellensatz} and~\ref{pb:cauchydav}; the binomial arithmetic of this chapter. The Catalan numbers return in Exercise~\ref{ex:catalan}.
\leads Erd\H{o}s and Heilbronn conjectured (a) in the $1960$s. Dias da Silva and Hamidoune proved it in $1994$ with exterior algebra. The polynomial argument above appeared a year later. The same idea, with a product $\prod_{i<j}(x_i-x_j)$, bounds sums of $h$ distinct elements. A direct check verifies (a) for every subset when $p\le11$.
\end{solution}

\begin{solution}{pb:sauer}
\idea Delete one point. The family shrinks only where two sets differ by that single point, and those pairs form a family of smaller VC-dimension.
\writeup (a) Let $\Phi(n,d)=\sum_{i=0}^d\binom ni$, with the conventions $\Phi(n,d)=2^n$ for $d\ge n$ and $\Phi(n,-1)=0$. Pascal's rule gives $\Phi(n,d)=\Phi(n-1,d)+\Phi(n-1,d-1)$. Induct on $n$. The case $n=0$ is immediate: the only subset of the empty set is the empty set.

Let $\mathcal F'=\{F\setminus\{n\}:F\in\mathcal F\}$ and let $\mathcal F''$ be the family of those $F\subseteq[n-1]$ for which both $F$ and $F\cup\{n\}$ lie in $\mathcal F$. The map $F\mapsto F\setminus\{n\}$ from $\mathcal F$ to $\mathcal F'$ is two-to-one exactly over $\mathcal F''$, and one-to-one elsewhere. So $|\mathcal F|=|\mathcal F'|+|\mathcal F''|$. A set shattered by $\mathcal F'$ is shattered by $\mathcal F$, so the VC-dimension of $\mathcal F'$ is at most $d$. If $\mathcal F''$ shatters $T\subseteq[n-1]$, then $\mathcal F$ shatters $T\cup\{n\}$, because both choices of whether to include $n$ are available. So the VC-dimension of $\mathcal F''$ is at most $d-1$. The inductive hypothesis gives $|\mathcal F|\le\Phi(n-1,d)+\Phi(n-1,d-1)=\Phi(n,d)$.

(b) Let the points be $x_1<\dots<x_n$. An interval cuts out either the empty set or a block $\{x_i,\dots,x_j\}$ with $i\le j$. The number of such sets is $1+n+\binom n2=\Phi(n,2)$. Any two of the points are shattered, by taking intervals that isolate one, the other, both, or neither. Among three points $x_i<x_j<x_k$, no interval cuts out exactly $\{x_i,x_k\}$. The VC-dimension is $2$, and the bound is attained.

\tools Pascal's rule, Proposition~\ref{prop:binom}.
\leads The lemma was found independently by Sauer, by Shelah, and by Vapnik and Chervonenkis around $1971$--$72$. In learning theory it is a reason why finite VC-dimension makes empirical averages converge uniformly. In discrete geometry it feeds the $\varepsilon$-net theorem of Haussler and Welzl.
\end{solution}

\begin{solution}{pb:bilinear}
\idea A bilinear exponential sum is small because Cauchy--Schwarz lets one variable run through the whole field. The resulting complete sum is a geometric series, and it vanishes unless the two frequencies agree. In two dimensions the same completion counts solutions of $u\cdot v=x$.
\writeup (a) Let $\sigma(a)=\sum_{b\in B}e(rab/p)$. The sum over $a\in A$ is complex, and Theorem~\ref{thm:cs} is the real form. Multiply that sum by a complex number of modulus $1$ so that it becomes real and nonnegative. Apply the theorem to those real parts and to a list of ones. A real part is at most the modulus, so
\[
\Bigl|\sum_{a\in A}\sigma(a)\Bigr|^2\le|A|\sum_{a\in A}|\sigma(a)|^2\le|A|\sum_{a\in\mathbb F_p}|\sigma(a)|^2.
\]
Expand the square on the right:
\[
\sum_{a\in\mathbb F_p}|\sigma(a)|^2=\sum_{b,b'\in B}\sum_{a\in\mathbb F_p}e\Bigl(\frac{ra(b-b')}{p}\Bigr).
\]
The inner sum is geometric. Let $\omega=e(r(b-b')/p)$. Then $\omega^p=e(r(b-b'))=1$, because $r(b-b')$ is an integer. The sum $\sum_{a=0}^{p-1}\omega^a$ equals $p$ if $\omega=1$, and equals $(\omega^p-1)/(\omega-1)=0$ if $\omega\ne1$. Now $\omega=1$ exactly when $p$ divides $r(b-b')$. The prime $p$ does not divide $r$, so this happens exactly when $b=b'$. There are $|B|$ such pairs, so the expanded sum equals $p|B|$. The factor $|A|$ from Cauchy--Schwarz then gives
\[
\Bigl|\sum_{a\in A}\sigma(a)\Bigr|^2\le p|A||B|.
\]
Both sides are nonnegative. The nonnegative square root is
\[
\Bigl|\sum_{a\in A}\sigma(a)\Bigr|\le\sqrt{p\,|A|\,|B|},
\]
which is the claim.

(b) Sum the geometric character over the residue $u\cdot v-x$:
\[
N(x)=\frac1p\sum_{r\in\mathbb F_p}\sum_{u\in U}\sum_{v\in V}e\Bigl(\frac{r(u\cdot v-x)}{p}\Bigr).
\]
The term $r=0$ equals $|U||V|/p$. Write $\hat U(w)=\sum_{u\in U}e(u\cdot w/p)$ for $w\in\mathbb F_p^2$. The remaining sum is $E=(1/p)\sum_{v\in V}\Sigma(v)$, where
\[
\Sigma(v)=\sum_{r\ne0}e\Bigl(-\frac{rx}{p}\Bigr)\hat U(rv).
\]
The same rotation as in (a), and then enlarging the sum from $V$ to all of $\mathbb F_p^2$, gives
\[
|E|^2\le\frac{|V|}{p^2}\sum_{v\in\mathbb F_p^2}|\Sigma(v)|^2.
\]
Expand the square and sum in $v$ first. The inner exponential sum over $v$ equals $p^2$ when $ru=r'u'$, and equals $0$ otherwise. So
\[
\sum_{v}|\Sigma(v)|^2=p^2\sum_{r,r'\ne0}e\Bigl(-\frac{(r-r')x}{p}\Bigr)\nu(r,r'),
\]
where $\nu(r,r')$ is the number of pairs $(u,u')\in U\times U$ with $ru=r'u'$. The diagonal $r=r'$ contributes $(p-1)|U|$. Off the diagonal, set $t=r/r'$, so $t\ne0,1$ and $\nu(r,r')$ depends only on $t$; call it $\mu(t)$. For fixed $t\ne1$, the sum over $r'\ne0$ of $e(-r'(t-1)x/p)$ equals $-1$, because $(t-1)x\not\equiv0$ and the complete sum over all residues is $0$, while the $r'=0$ term would have been $1$. The off-diagonal contribution is therefore $-\sum_{t\ne0,1}\mu(t)\le0$. Hence $|E|^2\le|V|(p-1)|U|\le p|U||V|$, because $p-1\le p$. Both sides of $|E|^2\le p|U||V|$ are nonnegative, so $|E|\le\sqrt{p\,|U|\,|V|}$. Since $N(x)-|U||V|/p=E$, this is the claim.

(c) Take $U=A\times C$ and $V=B\times D$. Then $N(x)$ counts the quadruples with $ab+cd=x$. Part (b) gives
\[
N(x)\ge\frac{|A||B||C||D|}{p}-\sqrt{p\,|A||B||C||D|},
\]
which is positive precisely when $|A||B||C||D|>p^3$. If $A=B=C=D$, the condition is $|A|^4>p^3$.

(d) For $1\le a,b\le m$ we have $0<ab/p\le m^2/p\le1/9$, so $0<2\pi ab/p\le2\pi/9$. Now $2\pi/9<\pi/3<\pi/2$, because $2/9<1/3<1/2$. Cosine is strictly decreasing on $[0,\pi/2]$, and $\cos(\pi/3)=1/2$, so $\cos(2\pi/9)>\cos(\pi/3)=1/2$. Each angle $2\pi ab/p$ is at most $2\pi/9$, so the real part of each summand is greater than $1/2$, and the real part of the sum is greater than $m^2/2$. With $m=\lfloor\sqrt p/3\rfloor$ we have $\sqrt p<3(m+1)$ whenever $m\ge1$, hence $\sqrt{p|A||B|}=\sqrt p\, m<(3m+3)m$. The sum is then larger than
\[
\frac{m}{6(m+1)}\sqrt{p|A||B|}.
\]
For $p\ge9$ one has $m\ge1$, and $m/(m+1)\ge1/2$, so the factor is at least $1/12$. The bound in (a) is therefore sharp up to a constant factor.

\tools Theorem~\ref{thm:cs}; the geometric sum from this chapter. Volume~5 develops the same completion as the inversion and Parseval theorem for the discrete Fourier transform.
\leads Bound (a) is a starting point of sum--product theory over $\mathbb F_p$. Part (b) is an incidence bound between points and lines in the plane $\mathbb F_p^2$.
\end{solution}

\section*{Sources and further reading}

\begin{itemize}\raggedright
\item R.~L.~Graham, D.~E.~Knuth and O.~Patashnik, \emph{Concrete Mathematics}, 2nd ed., Addison-Wesley, 1994, Chapters~2 and~5. Sums, binomial coefficients, and finite differences. Newton's formula is in Section~2.6. The convolution of Proposition~\ref{prop:vandermonde} is treated there under the name Chu--Vandermonde.
\item I.~Newton, \emph{Philosophi\ae\ naturalis principia mathematica}, London, 1687, Book~III, Lemma~V. The interpolation formula by finite differences.
\item A.~Engel, \emph{Problem-Solving Strategies}, Springer, 1998, Chapters~8 and~13. Induction, and a large stock of problems on sums.
\item A.~T.~Benjamin and J.~J.~Quinn, \emph{Proofs that Really Count: The Art of Combinatorial Proof}, Mathematical Association of America, 2003. Identities proved by double counting, including the Fibonacci identities of Exercises~\ref{ex:fibgcd} and~\ref{ex:fibsums}.
\item R.~B.~Nelsen, \emph{Proofs Without Words}, Mathematical Association of America, 1993. The pictures behind Figure~\ref{fig:gauss} and Exercise~\ref{ex:cubesgeo}.
\item R.~P.~Stanley, \emph{Catalan Numbers}, Cambridge University Press, 2015. Families of objects counted by the Catalan numbers, and the history: Euler's letter to Goldbach of 4 September 1751, and Catalan's memoir of 1838.
\item M.~Renault, \emph{Lost (and found) in translation: Andr\'e's actual method and its application to the generalized ballot problem}, Amer.\ Math.\ Monthly 115 (2008), 358--363. The reflection argument of Exercise~\ref{ex:catalan} is often credited to D.~Andr\'e (1887). Renault shows that Andr\'e's own argument was a different bijection.
\item A.~W.~F.~Edwards, \emph{Pascal's Arithmetical Triangle: The Story of a Mathematical Idea}, Johns Hopkins University Press, 2002. The triangle from Pingala, Jia Xian, al-Karaji and Khayyam through Pascal's \emph{Trait\'e}, written in 1654 and printed in 1665, including Pascal's statement of induction.
\item R.~Rashed, \emph{The Development of Arabic Mathematics: Between Arithmetic and Algebra}, Kluwer, 1994. Inductive arguments in al-Karaji and al-Samaw'al. Rashed reads the arguments as mathematical induction. The chapter uses the more cautious description: a step checked for the first exponents, together with the remark that the same step continues.
\item J.~Bernoulli, \emph{Positiones arithmeticae de seriebus infinitis}, Basel, 1689. The comparison of a geometric progression with the arithmetic progression that shares its first two terms.
\item Yang Hui, \emph{Xiangjie jiuzhang suanfa}, 1261. The binomial triangle printed there is attributed to Jia Xian. Zhu Shijie, \emph{Siyuan yujian}, 1303, prints the triangle again. Edwards, cited above, surveys both, together with al-Karaji, Khayyam and Pascal.
\item N.~L.~Rabinovitch, \emph{Rabbi Levi ben Gershon and the origins of mathematical induction}, Archive for History of Exact Sciences 6 (1970), 237--248. The \emph{Maasei Hoshev} of 1321.
\item Nicomachus of Gerasa, \emph{Introduction to Arithmetic}, translated by M.~L.~D'Ooge, Macmillan, 1926. Book~II, Chapter~20: the cubes as sums of consecutive odd numbers.
\item A.~Granville, \emph{Zaphod Beeblebrox's brain and the fifty-ninth row of Pascal's triangle}, Amer.\ Math.\ Monthly 99 (1992), 318--331. Pascal's triangle modulo primes and prime powers, the story behind Figure~\ref{fig:pascal}.
\item J.~W.~L.~Glaisher, \emph{On the residue of a binomial-theorem coefficient with respect to a prime modulus}, Quarterly Journal of Mathematics 30 (1899), 150--156. The count $2^{s(n)}$ of odd entries in row $n$.
\item A.~Granville and O.~Ramar\'e, \emph{Explicit bounds on exponential sums and the scarcity of squarefree binomial coefficients}, Mathematika 43 (1996), 73--107. The proof that $\binom{2n}n$ is never squarefree for $n>4$.
\item D.~E.~Knuth, \emph{Johann Faulhaber and sums of powers}, Mathematics of Computation 61 (1993), 277--294. What the \emph{Academia Algebrae} of 1631 contains.
\item S.~W.~Golomb, \emph{Checker boards and polyominoes}, Amer.\ Math.\ Monthly 61 (1954), 675--682. The tromino tiling.
\item L.~Theisinger, \emph{Bemerkung \"uber die harmonische Reihe}, Monatshefte f\"ur Mathematik und Physik 26 (1915), 132--134. An early proof that $H_n$ is not an integer, in the spirit of Exercise~\ref{ex:harmonic}.
\item L.~E.~Sigler, \emph{Fibonacci's Liber Abaci}, Springer, 2002. The rabbit problem is in Chapter~12 of the 1202 book. The name ``Fibonacci numbers'' was given by \'E.~Lucas in 1877.
\item B.~Hayes, \emph{Gauss's day of reckoning}, American Scientist 94 (2006), 200--205. The classroom story, from Sartorius von Waltershausen's \emph{Gauss zum Ged\"achtniss} (1856) onward, and how little of the familiar detail is in that first account.
\end{itemize}
